% These are "magic" comments understood by the "TeXStudio" editor
% !TeX encoding = UTF-8
% !TeX TS-program = lualatex
% !TeX TXS-program:bibliography = biber
% Document title: ETSETB TFG LaTeX Template
% Version: 5.1
% Author: 2023 Orestes Mas Casals
% License: ETSETB TFG LaTeX Template by Orestes Mas is marked with CC0 1.0 Universal

\documentclass[a4paper,twoside,12pt]{report}

% Load auxiliary packages for various tasks
% Document title: ETSETB TFG LaTeX Template
% Version: 5.0
% Author: 2023 Orestes Mas Casals
% License: ETSETB TFG LaTeX Template by Orestes Mas is marked with CC0 1.0 Universal

%%%% TOOLS %%%%
\usepackage{lipsum}
\usepackage{etoolbox}   % <<< Provides some macros to compare text
\usepackage{iftex}      % <<< Detect TeX engine used to compile the document

%%%% MATH SETUP %%%%
\usepackage{mathtools}

%%%% INTERNATIONAL SYSTEM NOTATION FOR UNITS IN NUMBERS %%%%
\usepackage{siunitx}

%%%% FONTS SETUP %%%%
\iftutex   % True if compiling with Unicode engines (XeTeX/LuaTeX)
  \usepackage{fontspec}
  \usepackage[math-style=ISO,
    warnings-off={mathtools-colon, mathtools-overbracket},
  ]{unicode-math}
  %  This package is only a wrapper for the two packages libertinus-type1 (pdfLaTeX) and libertinus-otf (LuaLaTeX/XeLaTeX).
  %  The Libertinus fonts are similiar to Libertine and Biolinum, but come with math symbols. 
  \usepackage{libertinus}
  % Use DejaVu Mono for monospaced font
  \usepackage[
    mono=true,serif=false,sans=false,math=false,
    TT={Scale=0.80,FakeStretch=1.0}]{dejavu-otf}
\else    % True if compiling with 8-bit engines (LaTeX/pdfLaTeX)
  \usepackage[T1]{fontenc}
  %  This package is only a wrapper for the two packages libertinus-type1 (pdfLaTeX) and libertinus-otf (LuaLaTeX/XeLaTeX).
  %  The Libertinus fonts are similiar to Libertine and Biolinum, but come with math symbols. 
  \usepackage{libertinus,libertinust1math}
  % Use DejaVu Mono for monospaced font
  \usepackage[scaled=0.80]{DejaVuSansMono}
  % Some declarations for the Gantt diagram example
  \DeclareUnicodeCharacter{1D49}{\textsuperscript{e}} % "ᵉ" character
  \DeclareUnicodeCharacter{02B3}{\textsuperscript{r}} % "ʳ" character
  \DeclareUnicodeCharacter{1D52}{\textsuperscript{o}} % "ᵒ" character
\fi
\usepackage{microtype}

%%%% PAGE GEOMETRY %%%%
\usepackage{geometry}

%%%% PAGE STYLE %%%%%
\usepackage{fancyhdr}

%%%% PARAGRAPH STYLE %%%%
% Blank line between paragraphs instead of first line indenting
\usepackage{parskip}

%%%% TITLES STYLE %%%%
\usepackage{titlesec}

%%%% COLOR DEFINITIONS %%%%
\usepackage[svgnames]{xcolor}

%%%% TABLES %%%%
%\usepackage{array}
\usepackage{tabularx,booktabs}

%%%% ACRONYMS %%%%
\usepackage{acro}
% class `abbrev': abbreviations:
\DeclareAcronym{EU}{
  short = EU,
  long  = European Union,
  tag = abbrev
}

\DeclareAcronym{ETSETB}{
  short = ETSETB,
  long  = Escola Tècnica Superior d'Enginyeria de Telecomunicació de Barcelona,
  tag = abbrev
}

\DeclareAcronym{GREELEC}{
  short = GREELEC,
  long  = Grau en Enginyeria Electrònica de Telecomunicació,
  tag = abbrev
}

\DeclareAcronym{GEF}{
  short = GEF,
  long  = Grau en Enginyeria Física,
  tag = abbrev
}

\DeclareAcronym{VQE}{
  short = VQE,
  long  = Variational Quantum Eigensolver,
  tag = abbrev
}

\DeclareAcronym{DQS}{
  short = DQS,
  long  = Digital Quantum Simulation,
  tag = abbrev
}

\DeclareAcronym{AQS}{
  short = AQS,
  long  = Analog Quantum Simulation,
  tag = abbrev
}

\DeclareAcronym{MPS}{
  short = MPS,
  long  = Matrix Product States,
  tag = abbrev
}

\DeclareAcronym{PEPS}{
  short = PEPS,
  long  = Projected Entangled Pair States,
  tag = abbrev
}

\DeclareAcronym{FCI}{
  short = FCI,
  long  = Full Configuration Interaction,
  tag = abbrev
}

\DeclareAcronym{CCSD}{
  short = CCSD,
  long  = Coupled Cluster with Singles and Doubles,
  tag = abbrev
}

\DeclareAcronym{ADAPT-VQE}{
  short = ADAPT-VQE,
  long  = Adaptive Variational Quantum Eigensolver,
  tag = abbrev
}

\DeclareAcronym{QNG}{
  short = QNG,
  long  = Quantum Natural Gradient,
  tag = abbrev
}

\DeclareAcronym{GD}{
  short = GD,
  long  = Gradient Descent,
  tag = abbrev
}

\DeclareAcronym{NMomentum}{
  short = NMomentum,
  long  = Nesterov Momentum,
  tag = abbrev
}

\DeclareAcronym{RMSProp}{
  short = RMSProp,
  long  = Root Mean Square Propagation,
  tag = abbrev
}

\DeclareAcronym{Adagrad}{
  short = Adagrad,
  long  = Adaptive Gradient Algorithm,
  tag = abbrev
}

\DeclareAcronym{Adam}{
  short = Adam,
  long  = Adaptive Moment Estimation,
  tag = abbrev
}

\DeclareAcronym{NISQ}{
  short = NISQ,
  long  = Noisy Intermediate-Scale Quantum,
  tag = abbrev
}

\DeclareAcronym{GPU}{
  short = GPU,
  long  = Graphics Processing Unit,
  tag = abbrev
}

\DeclareAcronym{CPU}{
  short = CPU,
  long  = Central Processing Unit,
  tag = abbrev
}

\DeclareAcronym{UCCSD}{
  short = UCCSD,
  long  = Unitary Coupled Cluster with Singles and Doubles,
  tag = abbrev
}

\DeclareAcronym{HF}{
  short = HF,
  long  = Hartree-Fock,
  tag = abbrev
}



%%%% GRAPHICS %%%%
\usepackage{graphicx}
\usepackage{caption}        %
\usepackage{tikz}           % Generic vector graphics language
\usepackage{pgfgantt}       % Gantt diagrams. See "pgfgantt" manual for reference

%%%% ELECTRONIC CIRCUITS %%%%
\usepackage{circuitikz}

%%%% PLOTTING %%%%
\usepackage{pgfplots}

%%%% COLOR BOXES %%%%
\usepackage{tcolorbox}

%%%% BIBLIOGRAPHY (using biblatex) %%%%
%\usepackage[section,numbib]{tocbibind}       % Makes bibliography appear in the Table of Contents
\usepackage[backend=biber,style=numeric]{biblatex}
\addbibresource{TFG.bib}

%%%% QUOTES %%%%
\usepackage{csquotes}

%%%% HYPERLINKS %%%%
\usepackage{hyperref}



% Document style definitions. GHANGE STYLES THERE.
% Document title: ETSETB TFG LaTeX Template
% Version: 5.0
% Author: 2023 Orestes Mas Casals
% License: ETSETB TFG LaTeX Template by Orestes Mas is marked with CC0 1.0 Universal

%%% COLORS
\definecolor[named]{ETSETBcopper}{RGB}{220,68,05}
\definecolor[named]{UPCblue}{RGB}{15,128,204}
\colorlet{DefaultColor}{SlateGray}
\colorlet{GRETSTcolor}{SteelBlue}
\colorlet{GREELECcolor}{SeaGreen}
\colorlet{GEFcolor}{SlateBlue}
% This color will change according to the Degree set.
\colorlet{DegreeColor}{ETSETBcopper}

%%% PAGE GEOMETRIES
% Geometry of cover page
\newgeometry{
  headheight=2cm,
  headsep=1cm,
  top=3cm,
  bottom=3cm,
  left=3.5cm,
  right=2.5cm,
  top=4.5cm,           % Allow space on the top for the logos
}
\savegeometry{cover}

% Geometry for main content pages
\newgeometry{%
  headheight=30pt,
  headsep=1cm,
  top=3.5cm,
  bottom=3.5cm,
  footskip=1.5cm,
  left=3cm,
  right=3cm,
  marginparwidth=20pt,
}
\savegeometry{main}

%%% Redefinition of various commands
%%% It must be done this way, because these styles are hardcoded into Book class
\makeatletter
\renewcommand{\tableofcontents}{%
  \chapter*{\contentsname%
    \@mkboth{\contentsname}{\contentsname}%
  }%
  \@starttoc{toc}%
}
\renewcommand{\listoffigures}{%
  \chapter*{\listfigurename%
    \@mkboth{\listfigurename}{\listfigurename}%
  }%
  \@starttoc{lof}%
}
\renewcommand{\listoftables}{%
  \chapter*{\listtablename%
    \@mkboth{\listtablename}{\listtablename}%
  }%
  \@starttoc{lot}%
}
\makeatother

%%% HEADER / FOOTER STYLES FOR VARIOUS DOCUMENTS PARTS

% "plain" page style is invoked always (unless changed by user) at the first page of each chapter
\fancypagestyle{plain}{
  \fancyhf{}
  \renewcommand\headrulewidth{0pt}
  \renewcommand\footrulewidth{0pt}
  \fancyfoot[C]{-- \thepage\ --}
}

% Front cover style
\fancypagestyle{frontcover}[plain]{%
  \fancyhf{}     % Empty current header and footer
  \fancyhead[L]{\includegraphics[height=1.75cm]{img/logos/Logo_UPC_ETSETB_tight.pdf}}
  \fancyhead[R]{\includegraphics[height=1.75cm]{img/logos/logo_telecos.png}}
  \renewcommand{\footrulewidth}{6pt}%
  \futurelet\TMPfootrule\def\footrule{{\color{gray!80}\TMPfootrule}}
}

% Main content style
\fancypagestyle{main}[plain]{%
  % Draw a line under the heading of every page EXCEPT on float pages or if they start with a figure
  \renewcommand\headrulewidth{\iftopfloat{0pt}{\iffloatpage{0pt}{1pt}}}
  \futurelet\TMPheadrule\def\headrule{{\color{gray!80}\TMPheadrule}}
  % 
  \renewcommand{\chaptermark}[1]{\markboth{\chaptername\ \thechapter{}. ##1}{}}
  \renewcommand{\sectionmark}[1]{\markright{\thesection\ - ##1}}
  \fancyhead[LE]{\scshape\leftmark}
  \fancyhead[RO]{\scshape\rightmark}%
}

%%% STYLE FOR CHAPTER TITLE
% Format of titles is:
% \titleformat
%   {command}
%   [shape]
%   {format}
%   {label}
%   {sep}
%   {before-code}
%   [after-code]
\titleformat
{\chapter}%
[display]        % The "display" shape puts the label in a separate paragraph
{\sffamily\normalsize\Huge\color{black}}% Format of the whole title
{\flushright\sffamily\normalsize\color{DegreeColor}\hrulefill\hspace{1ex}%
   \MakeUppercase{\chaptertitlename}\hspace{1ex}% 
   {\sffamily\fontsize{60}{60}\selectfont\thechapter}%
}% Label format - A line followed by a "Chapter" string and a large colored chapter number
{5 pt}% Vertical separator between the label and the title body
{\flushright\bfseries\LARGE}% Code preceding title body (title is LARGE bold and raggedleft)
[]% Code after the title body.

%%% STYLE FOR TABLES %%%
% Increase spacing between cell content and borders
\renewcommand{\arraystretch}{1.3}

%%% BIBLIOGRAPHY STYLE %%%
\ExecuteBibliographyOptions{  
  sorting=nyt,
}
% Enlarge separation between bibliography items
\setlength\bibitemsep{0.5\baselineskip}

%%% FIGURE's CAPTION STYLE %%%
\DeclareCaptionFont{blue}{\color{blue}}
\captionsetup{
  labelfont={blue,sf,bf},
  textfont=it,
}

%%% TIKZ/PGF DRAWING GLOBAL STYLES %%%
\usetikzlibrary{
  babel,  % DON'T EVER TRY TO REMOVE THAT IF YOU USE THE "babel" PACKAGE !!
  angles,calc,quotes,
  arrows.meta,positioning,shapes.misc,shapes.geometric,
  patterns,decorations.pathmorphing,decorations.markings,
}
\tikzset{
}

%%% PGFPLOTS SETTINGS %%%
\pgfplotsset{
  compat=1.18,       % Avoiding a warning
}

%%% GANTT DIAGRAMS GLOBAL STYLES %%%
\ganttset{
  progress label text = {\pgfmathprintnumber[precision=0, verbatim]{#1}\% \ifcase\doclanguage\or completa\or completa\else complete\fi}
}

%%% TCOLORBOX STYLES
\tcbuselibrary{skins,breakable,listingsutf8}

%%% HYPERREF OPTIONS %%%
\hypersetup{
  breaklinks=true,
  colorlinks=true,    % CHANGE TO "true" IF YOU WANT COLORED TEXT LINKS INSTEAD OF BEING SURROUNDED BY A BORDER.
%  hidelinks,          % UNCOMMENT IF YOU WANT TO REMOVE COLOR AND BORDER AROUND THE LINKS.
}


%%%%%%%%%%%%%%%%%%%%%%%%%%%%%%%%%%%%%%%%%%%%%%%%%%%%%%%%%%%%%%%%%%%%%%
%%%                                                                %%%
%%%              DADES DEL DOCUMENT / DOCUMENT DATA                %%%
%%%                                                                %%%
%%%%%%%%%%%%%%%%%%%%%%%%%%%%%%%%%%%%%%%%%%%%%%%%%%%%%%%%%%%%%%%%%%%%%%

%%% Please UNCOMMENT and/or provide suitable information to the macros below
%%% Si us plau, DESCOMENTEU i/o ompliu les macros següents amb la informació adient.

%%% Specify the document language
%%%   1: Catalan
%%%   2: Spanish
%%%   3: English
\newcommand{\doclanguage}{3}

% Set-up language and load language-related cover page macros
% Document title: ETSETB TFG LaTeX Template
% Version: 5.0
% Author: 2023 Orestes Mas Casals
% License: ETSETB TFG LaTeX Template by Orestes Mas is marked with CC0 1.0 Universal

%%%% TEMPLATE LANGUAGE MACROS %%%%

% Language selection based on the language variable defined above
\ifcase\doclanguage\or
  % Case 1
  \usepackage[catalan]{babel}
  \selectlanguage{catalan}%
  \setquotestyle{spanish}% El paquet "csquotes" no té estil pel català, però podem usar l'espanyol
  \newcommand\TFGdate{
    \ifcase\month\or%
      Gener\or Febrer\or Març\or Abril\or
      Maig\or Juny\or Juliol\or Agost\or
      Setembre\or Octubre\or Novembre\or Desembre\fi
      \space de~\number\year%
    }
\or % Case 2
  \usepackage[spanish,es-noquoting]{babel}
  \selectlanguage{spanish}%
  \newcommand\TFGdate{
    \ifcase\month\or%
      Enero\or Febrero\or Marzo\or Abril\or
      Mayo\or Junio\or Julio\or Agosto\or
      Septiembre\or Octubre\or Noviembre\or Diciembre\fi
      \space de~\number\year%
    }
\else % Case 3 (default)
  \usepackage[english]{babel}
  \selectlanguage{english}%
  \newcommand\TFGdate{
    \ifcase\month\or%
      January\or February\or March\or April\or
      May\or June\or July\or August\or
      September\or October\or November\or December\fi
      \space\number\year%
    }
\fi

\newcommand\degreename{
  \ifcase\doclanguage
  \or
    \{Nom del Grau\}%
  \or
    \{Nombre del Grado\}%
  \else
    \{Degree's Name\}%
  \fi
}

\newcommand\GEF{
  \ifcase\doclanguage
  \or
    Grau en Enginyeria Física
  \or
    Grado en Ingeniería Física
  \else
    Degree in Physics Engineering
  \fi
}

\newcommand\GRETST{
  \ifcase\doclanguage
  \or
    Grau en Enginyeria de Tecnologies i Serveis de Telecomunicació
  \or
    Grado en Ingeniería de Tecnologías y Servicios de Telecomunicación
  \else
    Degree in Telecommunication Technologies and Services Engineering
  \fi
}

\newcommand\GREELEC{
  \ifcase\doclanguage
  \or
    Grau en Enginyeria Electrònica de Telecomunicació
  \or
    Grado en Ingeniería Electrónica de Telecomunicación
  \else
    Degree in Electronics Engineering%  Should be "Degree in Telecommunication Electronics Engineering"
  \fi
}

\newcommand{\studentsname}{
  \ifcase\doclanguage
  \or
    \{Nom de l'Autor\}%
  \or
    \{Nombre del Autor\}%
  \else
    \{Author's Name\}%
  \fi
}

\newcommand\reporttitle{
  \ifcase\doclanguage
  \or
    \{Títol del Treball\}%
  \or
    \{Título del Trabajo\}%
  \else
    \{Report's title\}%
  \fi
}

\newcommand\FirstParagraph{
  \ifcase\doclanguage
  \or
    Treball de Fi de Grau\medskip\\
    presentat a l'Escola Tècnica Superior\medskip\\
    d'Enginyeria de Telecomunicació de Barcelona\medskip\\
    de la Universitat Politècnica de Catalunya\medskip\\
    per
  \or
    Trabajo de Final de Grado\medskip\\
    presentado en la Escola Tècnica Superior\medskip\\
    d'Enginyeria de Telecomunicació de Barcelona\medskip\\
    de la Universitat Politècnica de Catalunya\medskip\\
    por
  \else
    Bachelor's Degree Thesis\medskip\\
    Submitted to the Faculty at the\medskip\\
    Escola Tècnica Superior\medskip\\
    d'Enginyeria de Telecomunicació de Barcelona\medskip\\
    of the Universitat Politècnica de Catalunya\medskip\\
    by
  \fi
}

\newcommand\SecondParagraph{
  \ifcase\doclanguage
  \or
    En compliment parcial\medskip\\
    dels requisits per a l'obtenció del
  \or
    En cumplimiento parcial\medskip\\
    de los requisitos para la obtención del
  \else
    In partial fulfillment\medskip\\
    of the requirements for the
  \fi
}

\newcommand\advisor{
  \ifcase\doclanguage
  \or
    \textit{\{Nom del director del projecte\}}
  \or
    \textit{\{Nombre del director del proyecto\}}
  \else
    \textit{\{Advisor's Name\}}
  \fi
}

\newcommand\reviewer{
  \ifcase\doclanguage
  \or
    \textit{\{Nom del ponent (si s'escau)\}}
  \or
    \textit{\{Nombre del ponente (si aplica)\}}
  \else
    \textit{\{Reviewer's Name (if applicable)\}}
  \fi
}

\newcommand\ThesisTitle{\reporttitle}
\newcommand\AuthorName{\studentsname}
\newcommand\DegreeName{\degreename}
\newcommand\AdvisorName{\advisor}
\newcommand\ReviewerName{\reviewer}
\newcommand\documentDate{\TFGdate}
\newcommand\AdvisorLine{
  \ifcase\doclanguage
  \or
    Director/a:
  \or
    Director/a:
  \else
    Advisor:
  \fi \AdvisorName
}
\newcommand\ReviewerLine{
  \ifcase\doclanguage
  \or
    Ponent:
  \or
    Ponente:
  \else
    Reviewer:
  \fi \ReviewerName
}

% Define some macros for user-provided information
\newcommand\setTitle[1]{\renewcommand{\ThesisTitle}{#1}}
\newcommand\setAuthor[1]{\renewcommand{\AuthorName}{#1}}
\newcommand\setDegree[1]{%
  \renewcommand{\DegreeName}{#1}
  % This stablishes a different color for each Degree
  % It's only a proposal and proof-of-concept. By now it's disabled until approval
%  \ifdefstring{\DegreeName}{\GRETST}{
%    \colorlet{DegreeColor}{GRETSTcolor}
%  }{
%    \ifdefstring{\DegreeName}{\GREELEC}{
%      \colorlet{DegreeColor}{GREELECcolor}
%    }{
%      \ifdefstring{\DegreeName}{\GEF}{
%        \colorlet{DegreeColor}{GEFcolor}
%      }{
%        \colorlet{DegreeColor}{UPCblue}
%      }
%    }
%  }
}
\newcommand\setAdvisor[1]{\renewcommand{\AdvisorName}{#1}}
\newcommand\setReviewer[1]{\renewcommand{\ReviewerName}{#1}}
\newcommand\setDate[1]{\renewcommand{\documentDate}{#1}}



%%% Thesis Title / Títol del TFG
\setTitle{Quantum Simulation}

%%% Name of the Thesis Author / Nom de l'autor del TFG
\setAuthor{Albert López Escudero}

%%% Degree's name. Uncomment the "setDegree" macro and set its argument to one:
%%% Nom del Grau. Descomenteu la macro «setDegree» i passeu-li un d'aquests paràmetres:
%%%   \GEF     - Physics Engineering            <<<< DON'T USE THIS, as the template is not approved by the GEF faculty yet
%%%   \GREELEC - Electronics Engineering
%%%   \GRETST  - Telecommunications Engineering
\setDegree{\GREELEC}

%%% Name of the Thesis Advisor / Director
\setAdvisor{Vincenzo de Maio \& Ivona Bandric}

%%% Name of the Thesis Rewiewer / «Ponent»
\setReviewer{Javier Rodrigez Fonollosa}

%%% Document's date (\TFGdate by default)
\setDate{\TFGdate}


%%%%%%%%%%%%%%%%%%%%%%%%%%%%%%%%%%%%%%%%%%%%%%%%%%%%%%%%%%%%%%%%%%%%%%%%%%
%%%%%%%%%%%%%%%%%%%%% PREAMBLE ENDS - DOCUMENT BEGINS %%%%%%%%%%%%%%%%%%%%
%%%%%%%%%%%%%%%%%%%%%%%%%%%%%%%%%%%%%%%%%%%%%%%%%%%%%%%%%%%%%%%%%%%%%%%%%%

\begin{document}

%%%%%%%%%%%%%%%%%%%%%%%%% FRONT COVER %%%%%%%%%%%%%%%%%%%%%%%%%
\begin{titlepage}
  % Cover page style defined in config/styles.tex
  \loadgeometry{cover}
  \thispagestyle{frontcover}
  % Use a Sans Serif typeface throughout
  \sffamily
  
  % Title page background image
  \begin{tikzpicture}[remember picture,overlay]
    \node[anchor=center,opacity=0.3,xshift=5mm,yshift=-10mm,scale=8] (Campillo) at (current page.center) {\includegraphics{img/Campillo.png}};
    \filldraw[color=DegreeColor] (current page.south west) rectangle +(1.5cm,\paperheight);
    \filldraw[color=ETSETBcopper] (current page.south west) -- ++(0,0.140\paperheight-10pt) -- ++(0.75cm,20pt) --   ++(0.75cm,-20pt) -- ++(0,-0.140\paperheight+10pt) -- cycle;
    % This commented out sentence is work-in-progress
    % \node[anchor=south,at=(current page.south west),xshift=0.75cm,yshift=0.5cm] {\rotatebox{90}{\textcolor{white}{\Huge\bfseries \raisebox{-3pt}{ETSETB}~~~~\LARGE\DegreeName}}};
    \node[anchor=south,at=(current page.south west),xshift=0.75cm,yshift=0.5cm] {\rotatebox{90}{\textcolor{white}{\Huge\bfseries \raisebox{0pt}{ETSETB} -- \LARGE\DegreeName}}};
  \end{tikzpicture}
  
  \begin{center}
    \vspace*{2.5cm}
    
    {\Huge \ThesisTitle}\\
    
    \vspace{0.8cm}
    
    \color{black}\hrule height 2pt
    
    \vspace{1cm}
    
    {\large \FirstParagraph}\bigskip\\
  
    {\Large \AuthorName}\medskip\\
    
    \vspace{2.5cm}
    
    \SecondParagraph\medskip\\
    
    {\bfseries\scshape\DegreeName}
    
    % This is an elastic space
    \vfill
    
    \AdvisorLine\medskip\\
    
    %%% Comment out / delete the following line if there's no thesis reviewer
    \ReviewerLine\medskip\\
    
    \vspace{1cm}
    
    {Barcelona, \documentDate}
    
    \vspace{1cm}
  \end{center}
\end{titlepage}

% Restore default page geometry (modified by cover page code)
%\restoregeometry
\loadgeometry{main}

%%%%%%%%%%%%%%%%%%%%%%%%% PREAMBLE BEGINS %%%%%%%%%%%%%%%%%%%%%%%%%
\pagestyle{plain}


%%%% SUMMARY / RESUM %%%%
% DON'T CHANGE THIS LINE
\addcontentsline{toc}{section}{\ifcase\doclanguage\or Resum \or Resumen \else Summary\fi}

%%%% PLEASE REPLACE TEXTS WITH YOUR OWN CONTENT %%%%

%%% RESUM EN CATALÀ
\begin{center}
  \huge\bfseries\raggedright Resum~\hrulefill
\end{center}
  Cada exemplar del Treball de Fi de Grau (TFG) ha de contenir un Resum, que és un breu extracte del TFG. En termes d'estil, el Resum hauria de ser una versió reduïda del projecte: una introducció concisa, un compendi dels resultats i les principals conclusions o arguments presentats en el projecte. El Resum no ha de superar les 150 paraules i cal que estigui traduït al català, castellà i anglès.

%%% RESUMEN EN CASTELLANO
\begin{center}
  \huge\bfseries\raggedleft\vspace*{.5\baselineskip} \hrulefill ~Resumen
\end{center}
  Cada ejemplar del Trabajo de Fin de Grado (TFG) debe incluir un Resumenque es un breve extracto del TFG. En cuanto al estilo, el Resumen debería ser una versión reducida del proyecto: una introducción breve, un resumen de los resultados principales y las conclusiones o argumentos principales presentados en el proyecto. El Resumen no debe exceder las 150 palabras y debe estar traducido al catalán, castellano e inglés.

%%% ENGLISH SUMMARY
\begin{center}
  \huge\bfseries\raggedright\vspace*{.5\baselineskip} Summary~\hrulefill
\end{center}
  Each copy of the Bachelor's Thesis (TFG) must include a Summary, which is a concise abstract of the TFG. In terms of style, the Summary should be a condensed version of the project: a brief introduction, a summary of the main results, and the conclusions or key arguments presented in the project. The Summary should not exceed 150 words and must be translated to catalan, spanish and english.
                   % <<<<<< Please replace "summary.tex" entirely with your own content written in the required languages.


%%%% DEDICATION PAGE / PÀGINA DE DEDICATÒRIA %%%%
\clearpage
\thispagestyle{empty}             % No header/footer
\vspace*{4cm}                     % 
\hfill                            % Push the minipage to the right
\begin{minipage}{0.8\textwidth}   % You can adjust the minipage width
  \raggedleft\itshape             % Text in italics and flushed to the right
  \ifcase\doclanguage\or          % REPLACE FROM HERE (including the \ifcase structure) WITH YOUR OWN TEXT
    Podeu incloure una pàgina de Dedicatòries just abans de la pàgina d’Agraïments, però no és un requisit.\or
    Puede incluirse una página de Dedicatorias justo antes de la página de Agradecimientos, pero no es un requisito.\else
    A Dedication page may be included in your thesis just before the Acknowledgments page, but it is not a requirement.\fi
\end{minipage}



%%%% ACKNOWLEDGEMENTS / AGRAÏMENTS %%%%
\clearpage
\ifcase\doclanguage    % REPLACE FROM HERE (including the \ifcase structure) WITH YOUR OWN TEXT
\or
\chapter*{Agraïments}
\addcontentsline{toc}{section}{Agraïments}

És apropiat, però no obligatori, declarar l’extensió de l’ajuda aportada per persones de l'\textit{staff}, companys/companyes d’estudis, tècnics/ques o altres en la co\l.lecció de dades, disseny i construcció del prototip, l’anàlisi de dades, l’execució dels experiments i la preparació del projecte (incloent l’ajuda editorial). A més a més, és apropiat reconèixer la supervisió i la direcció donada pel tutor/a.
\or
\chapter*{Agradecimientos}
\addcontentsline{toc}{section}{Agradecimientos}

Es apropiado, aunque no obligatorio, declarar la extensión de la ayuda recibida por personas del staff, compañeros/as estudios, técnicos/as u otras en la recolección de datos, preparación del proyecto (incluyendo ayuda editorial). Además, es apropiado reconocer la supervisión, ayuda y dirección aportada por el tutor/a. 
\else
\chapter*{Acknowledgements}
\addcontentsline{toc}{section}{Acknowledgements}

It is appropriate, but not mandatory, to declare the extent to which assistance has been given by members of the staff, fellow students, technicians or others in the collection of materials and data, the design and construction of apparatus, the performance of experiments, the analysis of data, and the preparation of the thesis (including editorial help). In addition, it is appropriate to recognize the supervision and advice given by your advisor.
\fi



%%%% REVISION HISTORY TABLE %%%% <<<<<< Please read the instructions inside the "revision_history.tex" file.
\vspace*{0cm}

% As the Revision History has three LaTeX tables, in this particular case it's better to separate the translations in three different files.
% Please modify the one that suit your needs and this code will ignore the languages you're not working on.

\ifcase\doclanguage\or
  \vspace*{0cm}

\chapter*{Historial de revisió i aprovació}
\addcontentsline{toc}{section}{Historial de revisió i aprovació}

\begin{tabularx}{\textwidth}{|c|l|l|X|}
  \hline
    \textbf{Revisió} & \textbf{Data} & \textbf{Autor(s)} & \textbf{Descripció} \\
  \hline
    1.0 & dd/mm/yyyy & AME & Creació del document \\
  \hline
    1.1 & dd/mm/yyyy & AME, JPV & Correcció d'errors \\
  \hline
    2.0 & dd/mm/yyyy & AME, MLO & Versió revisada \\
  \hline
    4.0 & dd/mm/yyyy & AME & Versió final \\
  \hline
    & & & \\
  \hline
\end{tabularx}

\vspace{1cm}

\textbf{LLISTA DE DISTRIBUCIÓ DEL DOCUMENT}

\begin{tabularx}{\textwidth}{|l|X|}
  \hline
    \textbf{Rol} & \textbf{Cognom(s) i Nom} \\
  \hline
    [Estudiant] &  \\
  \hline
    [Director del projecte] &  \\
  \hline
    [Director 2 (si aplica)] &  \\
  \hline
\end{tabularx}

\vspace{1cm}

\begin{tabularx}{\textwidth}{|l|X||l|X|}
  \hline
    \multicolumn{2}{|l||}{\textbf{Escrit per:}}  &  \multicolumn{2}{l|}{\textbf{Revisat i aprovat per:}} \\
  \hline
    Data &  dd/mm/yyyy & Data &  dd/mm/yyyy \\
  \hline
    Nom &  Xxxxxxx Yyyyyyy & Nom &  Xxxxxxx Yyyyyyy \\
  \hline
    Rol & Autor del projecte & Rol & Director del projecte \\
  \hline
\end{tabularx}


\or
  \vspace*{0cm}

\chapter*{Historial de revisión y aprobación}
\addcontentsline{toc}{section}{Historial de revisión y aprobación}

\begin{tabularx}{\textwidth}{|c|l|l|X|}
  \hline
    \textbf{Revisión} & \textbf{Fecha} & \textbf{Autor(es)} & \textbf{Descripción} \\
  \hline
    1.0 & dd/mm/yyyy & AME & Creación del documento \\
  \hline
    1.1 & dd/mm/yyyy & AME, JPV & Corrección de errores \\
  \hline
    2.0 & dd/mm/yyyy & AME, MLO & Versión revisada \\
  \hline
    4.0 & dd/mm/yyyy & AME & Versión final \\
  \hline
    & & & \\
  \hline
\end{tabularx}

\vspace{1cm}

\textbf{LISTA DE DISTRIBUCIÓN DEL DOCUMENTO}

\begin{tabularx}{\textwidth}{|l|X|}
  \hline
    \textbf{Rol} & \textbf{Apellido(s) y Nombre} \\
  \hline
    [Estudiante] &  \\
  \hline
    [Director del proyecto] &  \\
  \hline
    [Director 2 (sólo si aplica)] &  \\
  \hline
\end{tabularx}

\vspace{1cm}

\begin{tabularx}{\textwidth}{|l|X||l|X|}
  \hline
    \multicolumn{2}{|l||}{\textbf{Escrito por:}}  &  \multicolumn{2}{l|}{\textbf{Revisado y aprobado por:}} \\
  \hline
    Fecha &  dd/mm/yyyy & Fecha &  dd/mm/yyyy \\
  \hline
    Nombre &  Albert López Escudero & Nombre &  Xxxxxxx Yyyyyyy \\
  \hline
    Rol & Autor del proyecto & Rol & Director del proyecto \\
  \hline
\end{tabularx}


\else
  \vspace*{0cm}

\chapter*{Revision history and approval record}
\addcontentsline{toc}{section}{Revision history and approval record}
  
\begin{tabularx}{\textwidth}{|c|l|l|X|}
  \hline
  \textbf{Revision} & \textbf{Date} & \textbf{Author(s)} & \textbf{Description} \\
  \hline
  1.0 & dd/mm/yyyy & AME & Document creation \\
  \hline
  1.1 & dd/mm/yyyy & AME, JPV & Error correction \\
  \hline
  2.0 & dd/mm/yyyy & AME, MLO & Revised after review \\
  \hline
  4.0 & dd/mm/yyyy & AME & Final version \\
  \hline
  & & & \\
  \hline
\end{tabularx}

\vspace{1cm}

\textbf{DOCUMENT DISTRIBUTION LIST}

\begin{tabularx}{\textwidth}{|l|X|}
  \hline
    \textbf{Role} & \textbf{Surname(s) and Name} \\
  \hline
    [Student] &  \\
  \hline
    [Project Supervisor 1] &  \\
  \hline
    [Project Supervisor 2 (if applicable)] &  \\
  \hline
\end{tabularx}

\vspace{1cm}

\begin{tabularx}{\textwidth}{|l|X||l|X|}
  \hline
    \multicolumn{2}{|l||}{\textbf{Written by:}}  &  \multicolumn{2}{l|}{\textbf{Reviewed and approved by:}} \\
  \hline
    Date &  dd/mm/yyyy & Date &  dd/mm/yyyy \\
  \hline
    Name &  Xxxxxxx Yyyyyyy & Name &  Xxxxxxx Yyyyyyy \\
  \hline
    Position & Project Author & Position & Project Supervisor \\
  \hline
\end{tabularx}

\fi


%%%% INDEX %%%%
\clearpage
\addcontentsline{toc}{section}{\ifcase\doclanguage\or Índex\or Índice\else Contents\fi}
\tableofcontents


%%%% LISTS %%%%
\listoffigures 
\addcontentsline{toc}{section}{\ifcase\doclanguage\or Índex de figures \or Lista de figuras \else List of figures\fi}
\listoftables
\addcontentsline{toc}{section}{\ifcase\doclanguage\or Índex de taules \or Índice de cuadros \else List of tables\fi}
%\lstlistoflistings


%%%% ACRONYMS %%%% <<<<<< PLEASE DEFINE YOUR ACRONYMS IN "acronyms.tex"
\clearpage
\newcommand\AbbrvName{\ifcase\doclanguage\or Sigles i acrònims\or Siglas y acrónimos\else Abbreviations\fi}
\addcontentsline{toc}{section}{\AbbrvName}
\printacronyms[
  heading=chapter*,
  display=all,
  include=abbrev,
  name=\AbbrvName,
]


%%%%%%%%%%%%%%%%%%%%%%%%%%%%%%%%%%%%%%%%%%%%%%%%%%%%%%%%%%%%%%%%%%%%
%%%%%%%%%%%%%%%%%%%%%%%%% MAIN PART BEGINS %%%%%%%%%%%%%%%%%%%%%%%%%
%%%%%%%%%%%%%%%%%%%%%%%%%%%%%%%%%%%%%%%%%%%%%%%%%%%%%%%%%%%%%%%%%%%%

% Starts normal numbering for the rest of the document
%\mainmatter  (valid/useful only if using the "book" class instead of "report" in the very first line)
\pagestyle{main}


% Introduction. Please replace "introduction.tex" entirely with your own content written in the desired language.
%%%% PLEASE REPLACE ENTIRELY WITH YOUR OWN CONTENT %%%%

\chapter{Introduction}

In recent decades, computing has undergone a significant transformation, reaching physical limits in the improvement of classical hardware. These challenges have driven research in quantum computing, an emerging technology with the potential to revolutionize various fields due to its ability to perform calculations exponentially more efficiently in certain problems. Among these applications, quantum molecular simulation stands out as a promising area, enabling the study of complex chemical systems that are intractable for classical methods.

The main objective of this project is to develop a quantum simulator capable of efficiently modeling molecules using advanced quantum algorithms, with a specific focus on implementing the Variational Quantum Eigensolver (VQE) algorithm. This algorithm combines quantum circuits for state preparation and classical optimizers that adjust parameters to minimize the system's energy. Additionally, the project introduces innovations in the design of adaptive quantum circuits, the integration of simultaneous optimizations of nuclear and electronic coordinates, and visualization tools for analyzing the optimization process.

\section{Work goals}
The primary goal of this project is to contribute a novel methodology to the field of molecular simulations by:

\begin{itemize} 
  \item Developing a quantum simulation program capable of efficiently modeling molecules using quantum algorithms. 
  \item Introducing innovative code that implements adaptive circuit construction and advanced optimization strategies. 
  \item Demonstrating the effectiveness of this approach by simulating selected molecular systems and comparing the results with classical computational methods.
  \item Enhancing the integration between quantum and classical computations to optimize both electronic and nuclear degrees of freedom. 
\end{itemize}

\section{Requirements and specifications}

To achieve these objectives, the following requirements and specifications have been established:

\begin{itemize} 
  \item \textbf{Quantum Computing Frameworks}: Utilize quantum computing libraries such as PennyLane and JAX for quantum circuit simulation and automatic differentiation. 
  \item \textbf{Algorithm Implementation}: Implement the VQE algorithm with adaptive circuit construction, allowing the selection of the most significant excitations based on energy gradients. 
  \item \textbf{Optimization Techniques}: Employ efficient optimizers, including gradient-based methods like gradient descent and advanced techniques like the Quantum Natural Gradient optimizer. 
  \item \textbf{Hamiltonian Construction}: Accurately build molecular Hamiltonians for various molecular geometries and ensure compatibility with standard quantum chemistry basis sets. 
  \item \textbf{Visualization Tools}: Develop modules for visualizing energy convergence, parameter evolution, and molecular geometries throughout the optimization process. 
  \item \textbf{Computational Resources}: Ensure the code is optimized for performance, making effective use of computational resources and supporting parallel execution where possible. 
  \item \textbf{Extensibility}: Design the codebase to be modular and extensible, allowing future enhancements and adaptation to other molecular systems or quantum algorithms. 
\end{itemize}

\section{Methods and procedures}

This project employs innovative methods and procedures to enhance molecular simulations using quantum computing. The implementation reflects a cohesive integration of adaptive circuit construction, advanced optimization strategies, and a seamless blend of quantum and classical computation, as demonstrated by the project's modular codebase.

\subsection*{Adaptive Circuit Construction}

Traditional VQE implementations rely on fixed ansätze, which may fail to capture the complexities of molecular systems efficiently. In this project, adaptive circuit construction is implemented in the ansatz\_preparer.py module. The quantum circuit is dynamically constructed by selecting excitations (operators) from a predefined pool, based on their contribution to lowering the system's energy. This selection is guided by energy gradient computations provided by optimizer.py and hamiltonian\_builder.py. This approach ensures computational efficiency by including only the most impactful excitations, significantly accelerating convergence to the ground state energy.

\subsection*{Advanced Optimization Techniques}

Optimizing quantum circuit parameters is critical for the success of VQE. The optimizer.py and opt\_mol.py modules incorporate both gradien\-based and gradient\-free optimization methods. Techniques such as the Gradient Descent Optimizer and Quantum Natural Gradient Optimizer are employed to address the geometric structure of the parameter space, ensuring faster and more stable convergence. Additionally, the optimization process extends to nuclear coordinates, implemented in opt\_mol.py, allowing for simultaneous refinement of electronic structures and molecular geometries. This hybrid optimization loop is a key innovation, enabling the algorithm to explore the coupled energy landscape effectively.

\subsection*{Integration of Quantum and Classical Computation}

The seamless integration of quantum and classical computations is achieved using frameworks like PennyLane. This integration, implemented across hamiltonian\_builder.py, ansatz\_preparer.py, and optimizer.py, allows for efficient computation of energy gradients with respect to quantum circuit parameters and nuclear coordinates. These automatic differentiation tools facilitate robust optimization and enable handling of complex molecular systems that are challenging for classical methods.

\subsection*{Innovations in Code Implementation}

Key innovations in our code include:

\begin{itemize} 
  \item \textbf{Dynamic Operator Selection}: A procedure to compute energy gradients for each operator in the pool and select the one with the highest impact, thus adaptively constructing the quantum circuit. 
  \item \textbf{Hybrid Optimization Loop}: An iterative loop that updates both quantum circuit parameters and nuclear positions, enhancing the ability to find the global minimum energy configuration. 
  \item \textbf{Efficient Gradient Computation}: Advanced numerical methods in optimizer.py compute gradients for nuclear coordinates, enabling efficient geometry optimization.
  \item \textbf{Visualization and Analysis Tools}: The \textit{visualizer.py} module provides comprehensive visualization tools for analyzing energy evolution and 3D representations of molecular geometries.
  \item \textbf{Modularity and Extensibility}: The modular code structure, described in \textit{main.py} and related modules, facilitates easy adaptation to different molecules, basis sets, and quantum devices, supporting future research. 
\end{itemize}

\subsection*{Utilization of Existing Frameworks and Contribution to the Field}

The project builds on established frameworks such as PennyLane and leverages existing algorithms like the VQE. The hamiltonian\_builder.py module utilizes PennyLane's tools for molecular Hamiltonian construction, while the adaptive circuit construction and hybrid optimization approaches extend these capabilities significantly. By introducing novel methodologies and improving computational efficiency, this project makes a meaningful contribution to the fields of quantum chemistry and quantum computing.

\section{Work plan}

  \label{sec:workplan}

Normally the figures and tables are put in \verb|\figure| and \verb|\table| environments, that can float freely in the document. You can identify each float with a \verb|\label|

\begin{figure}[H]
  \centering
  \includegraphics[width=1\textwidth]{img/Gantt_diagram.png}
  \caption{Project's Gantt diagram{\footnotesize{Gantt diagram of the project. For more information read the manual \cite{skalagantt} of Skala.}.}}
  \label{fig:gantt_diagram}
\end{figure}




% State of the art chapter. Please replace "state_of_art.tex" entirely with your own content written in the desired language.
\chapter[State of the Art]{State of the Art}

In this project, we focus on the operation of quantum simulators applied to molecular simulation, an area where quantum computing offers significant advantages. The main objective is to develop a program capable of simulating different molecules in a simple and efficient manner.

To this end, we will review the basic concepts of quantum mechanics that are essential to understand the fundamentals and potential of quantum computing, and explore the techniques of quantum simulation that allow us to harness quantum computing capabilities even with current technological limitations.

\section{Quantum Computing}

It is essential to understand the difference between bits in classical computing and qubits in quantum computing to delve into this new technological paradigm.

In classical computing, the basic unit of information is the \textbf{bit}, which can take the value of 0 or 1. These bits are the foundation upon which conventional computers operate, processing information through combinations of these binary states.

In contrast, quantum computing uses the \textbf{qubit} or quantum bit as its basic unit. Unlike the classical bit, a qubit can exist in a superposition of states, meaning it can simultaneously represent the values 0 and 1 thanks to the principle of superposition in quantum mechanics. This property, along with phenomena such as quantum entanglement and interference, allows quantum computers to process information exponentially more efficiently for certain problems.

\begin{figure}[H]
    \centering
    \includegraphics[width=0.8\textwidth]{img/bit_vs_qbit.png}
    \caption{Comparison between the growth of classical and quantum computational capacity.\cite{mooreslaw}}
    \label{fig:bit_vs_qubit}
\end{figure}

This graph shows us the boundaries that quantum computing allows us to reach compared to classical computing. It could achieve and dramatically surpass computational power.

Understanding how qubits operate and their differences from classical bits is essential to understand the potential of quantum computing. For this reason, the fundamental concepts of quantum computing are presented below.


\subsection{Qubit}

The \textbf{qubit} is the basic unit of information in quantum computing. While the classical bit can only be in one of two states (0 or 1), a qubit can be in a superposition of both states simultaneously. This is due to the principle of quantum superposition, one of the fundamental characteristics of quantum mechanics. It is worth noting that the states of a system are usually described by vectors, which reside in a Hilbert space.

For qubits, the Hilbert space is 2-dimensional, and the following values can be taken:

\begin{itemize}
    \item \(\ket{0}\): Represents the basis state 0.
    \[
    \ket{0} = \begin{pmatrix} 1 \\ 0 \end{pmatrix}
    \]

    \item \(\ket{1}\): Represents the basis state 1.
    \[
    \ket{1} = \begin{pmatrix} 0 \\ 1 \end{pmatrix}
    \]
\end{itemize}

Mathematically, a qubit is represented as a linear combination of the basis states $\ket{0}$ and $\ket{1}$:

\[
\ket{\psi} = \alpha\ket{0} + \beta\ket{1}
\]

where $\alpha$ and $\beta$ are complex numbers that satisfy the normalization condition $|\alpha|^2 + |\beta|^2 = 1$. These coefficients indicate the probability amplitudes of finding the qubit in the states $\ket{0}$ or $\ket{1}$ upon measurement.\cite{wikipedia_qubit}

\subsection{Quantum Entanglement}
An important distinguishing feature between qubits and classical bits is that multiple qubits can exhibit \textbf{quantum entanglement}; the qubit itself is an exhibition of quantum entanglement. In this case, quantum entanglement is a local or nonlocal property of two or more qubits that allows a set of qubits to express higher correlation than is possible in classical systems.

The simplest system to display quantum entanglement is the system of two qubits. Consider, for example, two entangled qubits in the \(\ket{\Phi^+}\) Bell state:
\[
\ket{\Phi^+} = \frac{1}{\sqrt{2}} \big( \ket{00} + \ket{11} \big).
\]
There are equal probabilities of measuring either product state \(\ket{00}\) or \(\ket{11}\), this is called \textbf{equal superposition}:
\[
\left| \frac{1}{\sqrt{2}} \right|^2 = \frac{1}{2}.
\]
In other words, there is no way to tell if the first and the second qubit has value 0 or 1.

To better understand quantum entanglement, let's explain an example. Imagine these qubits are separated, with one qubit given to Alice and the other to Bob. Alice performs a measurement, and the result she obtains has equal probabilities for both \(\ket{0}\) and \(\ket{1}\). If Bob measures the value of his qubit at the same time, the result must be exactly the same due to quantum entanglement.

If Alice's measurement had been \(\ket{0}\), then Bob's result would also have been \(\ket{0}\), because \(\ket{00}\) is the only state where Alice's qubit is \(\ket{0}\). In short, for these two entangled qubits, whatever Alice measures, so does Bob, demonstrating perfect correlation in any basis, regardless of the distance between them. This is a surprising phenomenon that cannot be explained by classical physics.\cite{pirandola2015advances}


\subsection{Quantum Superposition}

\textbf{Quantum superposition} allows a quantum system to exist in multiple states simultaneously until a measurement is performed. This characteristic is key to the functioning of quantum computers, as it enables processing a large amount of inputs simultaneously.

A qubit, by itself, is not much more powerful than a conventional bit. What makes it truly powerful is that it can place the quantum information it holds into a state of superposition, which represents a combination of all possible configurations of the qubit. Groups of qubits in superposition can create complex, multidimensional computational spaces, allowing complex problems to be represented in new ways within these spaces.\cite{ibm_quantum_computing}

Superposition is especially useful in simulating complex molecular systems. Quantum computers can naturally model the superpositions of electronic states in molecules, which is crucial for studying chemical reactions and molecular properties that are difficult to address with classical methods due to the exponential growth of computational resources required.

\subsection{Quantum Decoherence}

\textbf{Quantum decoherence} is one of the main challenges in quantum computing. Such machines are expected to rely heavily on the undisturbed evolution of quantum coherence. Decoherence causes the system to lose its quantumness, which invalidates the superposition principle and turns quantum to classical behavior. They must handle the decoherence in order to maintain quantum computation performance.\cite{optimizing_decoherence}

Qubits are extremely sensitive to external disturbances, such as electromagnetic fluctuations, vibrations, and temperature changes. These interactions can cause quantum states to mix with those of the environment, leading to a loss of coherence that is irreversible and degrades the stored quantum information.\cite{decoherence}

To mitigate the effects of decoherence, various strategies are implemented:

\begin{itemize}
    \item \textbf{System Isolation}: Designing physical systems that minimize unwanted interactions with the environment, using materials and techniques that protect qubits from external disturbances.
    \item \textbf{Quantum Error Correction}: Implementing error correction codes that allow detecting and correcting errors without directly measuring the qubit's state, thereby preserving quantum information.
    \item \textbf{Dynamic Control}: Applying techniques such as pulse refocusing and dynamic pulse sequences that actively compensate for disturbances and extend the coherence time of qubits.
\end{itemize}

Controlling and mitigating decoherence are essential for the advancement of quantum computing and its application in areas like molecular simulation, where the precision of calculations is fundamental.


\section{Quantum Simulation}

Quantum simulation has become a vital tool for investigating intricate quantum systems that are beyond the reach of classical computational methods. Building on Richard Feynman’s idea that a computer operating on quantum principles could exceed the capabilities of classical machines for modeling quantum behavior, the field has seen rapid advancements. Today, it spans both digital and analog approaches and finds applications in numerous scientific domains.

There are mainly two approaches in quantum simulation: \textbf{Digital Quantum Simulation (DQS)} and \textbf{Analog Quantum Simulation (AQS)}. DQS employs quantum circuit models to study the dynamics of the system, thus emulating, with the help of qubits, the behavior of the system. This approach is universal, as it can, in principle, simulate any quantum system, although not always efficiently. On the other hand, AQS uses an accessible quantum system to emulate the Hamiltonian of the system, thereby allowing us to learn about the properties of the system we want to study. This method is particularly useful when a qualitative representation is required rather than high precision.\cite{AQS}

In addition to these approaches, there are algorithms inspired by quantum information theory that facilitate the classical simulation of quantum systems using the internal symmetries of particles. Techniques such as \textbf{Matrix Product States (MPS)} and \textbf{Projected Entangled Pair States (PEPS)} allow representing particle systems on classical computers more efficiently than standard classical methods, optimizing the calculation of properties of complex quantum systems.
\cite{mps_peps_tensor_networks}

The applications of the quantum simulation systems are an exciting area of scientific practice. In condensed matter physics, it allows the study of models such as the Hubbard model and quantum phase transitions, fundamental for understanding phenomena like superconductivity, it facilitates the calculation of molecular energies and complex chemical reactions, it emulates particles in high-energy fields and cosmological phenomena. Furthermore, quantum simulation promises to have applications in the study of many problems; high-energy physics, atomic physics, quantum chemistry, condensed-matter physics and cosmology. This could be implemented using quantum computers, but also with
simpler, analog devices that require less control and are, therefore, easier to construct.\cite{QuantumSimulation}

But there are still challenges to deal with. It’s tough to control quantum simulator systems properly. Plus, errors and decoherence can degrade or adversely affect the results. The number of qubits and quantum gates we need depends on how big and complex the system is. We think you need around 40 to 100 qubits to beat classical computers on some problems. Despite these hurdles, tech is advancing, and we’re getting better at quantum simulation. This could really change how we do research in science and help us learn more about the quantum world.

\section{Hamiltonian}

The Hamiltonian is a fundamental concept originating from classical mechanics, who developed a reformulation of Newtonian mechanics. Hamiltonian mechanics is a reformulation of classical mechanics that provides powerful tools for studying the dynamics of systems. 

The Hamiltonian of a system represents the total energy of the system, expressed in terms of generalized coordinates and momenta, and is given by the sum of the kinetic and potential energies. The \textbf{Hamiltonian operator} plays a central role in describing the energy and time evolution of quantum systems, takes different forms and can be simplified, in some cases. It represents the total energy of the system, including kinetic and potential energies, being essential for formulating the Schrödinger equation.\cite{wiki:hamiltonian_quantum_mechanics}

\subsection{Mathematical Definition}

The Hamiltonian operator, typically denoted as \( \hat{H} \), is a self-adjoint operator acting on the Hilbert space associated with the quantum system. For a single particle in one dimension, the Hamiltonian is expressed as:

\[
\hat{H} = \hat{T} + \hat{V}
\]

where:

\begin{itemize}
    \item \( \hat{T} \) is the kinetic energy operator.
    \item \( \hat{V} \) is the potential energy operator.
\end{itemize}

In terms of the position \( \hat{x} \) and momentum \( \hat{p} \) operators, these are defined as:

\[
\hat{T} = \frac{\hat{p}^2}{2m} = -\frac{\hbar^2}{2m} \frac{d^2}{dx^2}
\]

\[
\hat{V} = V(\hat{x})
\]

Here, \( m \) is the mass of the particle, \( \hbar \) is the reduced Planck constant, and \( V(\hat{x}) \) is the potential energy function depending on position.

Finally, it can be shown that the expectation value of the Hamiltonian, which represents the energy expectation value, is always greater than or equal to the system's minimum potential:\cite{wiki:hamiltonian_quantum_mechanics}

Computing the expectation value of the kinetic energy, we have:
\[
\begin{aligned}
T &= -\frac{\hbar^2}{2m} \int_{-\infty}^{+\infty} \psi^* \left( \frac{d^2 \psi}{dx^2} \right) \, dx \\ 
  &= -\frac{\hbar^2}{2m} \left( \left[ \psi'(x) \psi^*(x) \right]_{-\infty}^{+\infty} 
      - \int_{-\infty}^{+\infty} \left( \frac{d\psi}{dx} \right) 
      \left( \frac{d\psi}{dx} \right)^* \, dx \right) \\ 
  &= \frac{\hbar^2}{2m} \int_{-\infty}^{+\infty} \left| \frac{d\psi}{dx} \right|^2 \, dx \geq 0.
\end{aligned}
\]
Hence the expectation value of the kinetic energy is always non\-negative. We can calculate the expectation value of the total enegy for normalized wavefunction:

\[
E = T + \langle V(x) \rangle = T + \int_{-\infty}^{+\infty} V(x) |\psi(x)|^2 \, dx 
\geq V_{\text{min}}(x) \int_{-\infty}^{+\infty} |\psi(x)|^2 \, dx \geq V_{\text{min}}(x).
\]


\subsection{Role in the Schrödinger Equation}

The Hamiltonian is central to the Schrödinger equation, which describes how the quantum state of a system evolves over time. The time-dependent Schrödinger equation is expressed as:

\[
i\hbar \frac{\partial}{\partial t} |\psi(t)\rangle = \hat{H} |\psi(t)\rangle
\]

where \( |\psi(t)\rangle \) is the state vector of the system at time \( t \). For time-independent systems, the Schrödinger equation reduces to the eigenvalue equation:

\[
\hat{H} |\psi\rangle = E |\psi\rangle
\]

Here, \( E \) represents the energy level or eigenvalues of the state $\psi$ , this equation is known as the \textit{time-intependent Schrödinger equation}.\cite{tao_schrodinger}

\subsection{Hamiltonian in Multi-Particle Systems}

We can extend the concept of the Hamiltonian to systems with N particles.\cite{wiki:hamiltonian_quantum_mechanics} The Hamiltonian is expressed as:

\[
\hat{H} = \sum_{n=1}^{N} \hat{T}_{n} + \hat{V}
\]

where:
\begin{itemize}
    \item \(\hat{V}\): is the potential energy function, which depends on the spatial configuration of the system and time. A particular set of spatial positions at a given instant defines a configuration.
    \item \(\hat{T}_{n}\): is the kinetic energy operator for particle \(n\), given by:
    \[
    \hat{T}_{n} = \frac{\mathbf{\hat{p}}_{n} \cdot \mathbf{\hat{p}}_{n}}{2m_{n}} = -\frac{\hbar^2}{2m_{n}} \nabla_{n}^2
    \]
\end{itemize}

Combining these yields, the Schrödinger Hamiltonian for an N number of particles system:
\[
\begin{aligned}
\hat{H} &= \sum_{n=1}^{N} \hat{T}_{n} + \hat{V} \\[6pt]
        &= \sum_{n=1}^{N} \frac{\mathbf{\hat{p}}_{n} \cdot \mathbf{\hat{p}}_{n}}{2m_{n}} + V(\mathbf{r}_{1}, \mathbf{r}_{2}, \ldots, \mathbf{r}_{N}, t) \\[6pt]
        &= -\frac{\hbar^2}{2} \sum_{n=1}^{N} \frac{1}{m_{n}} \nabla_{n}^2 + V(\mathbf{r}_{1}, \mathbf{r}_{2}, \ldots, \mathbf{r}_{N}, t)
\end{aligned}
\]
\subsection{The second quantization}

Finding an exact solution of the  Schrödinger equation is equivalent to solving the Full Configuration Interaction (FCI) functions, where the wave function of a molecule is represented as a sum of all possible Slater determinants that can be formed using a specific set of molecular orbitals. As the number of electrons and orbitals increases, the number of Slater determinants grows exponentially, making the FCI method computationally infeasible for large systems. Therefore, it is recommended to use a quantum computer to solve this problem. For this reason, we need to transform the Hamiltonian to the second quantized operators' form.

For more than one particle, the second\-quantized Hamiltonian can be written under the form:
\[
H =
\sum_{p,q} h_{pq} a_p^\dagger a_q +
\sum_{p,q,r,s} h_{pqrs} a_p^\dagger a_q^\dagger a_r a_s,
\]
with \(a^\dagger\) and \(a\) being the electron creation and annihilation operators. The first term thus represents the transitions of single electrons between orbitals, while the second term corresponds to the interactions between pairs of electrons\cite{SecondQuantization}.

\subsection{Importance in Quantum Simulations}

In quantum simulations, especially in algorithms like the Variational Quantum Eigensolver (VQE), the Hamiltonian is decomposed into a sum of simpler terms, often expressed in terms of Pauli operators. This decomposition facilitates implementation on quantum circuits and allows estimating the system's energy through measurements on qubits.

Understanding the structure and properties of the Hamiltonian is essential for modeling and simulating quantum systems, as it determines the possible energies and dynamics of the system under study.

\section{VQE: Variational Quantum Eigensolver}
\label{sec:vqe}

Variational quantum eigensolvers (VQEs) are among the most promising candidates for achieving useful computations in chemistry on near-term quantum computers. They operate by varying parameters and optimizing them to achieve the lowest energy, this can be achived by using variational principle from quantum mechanics. 

At their core, they prepare a quantum state, determined by a set of classical parameters that are specified by an \textit{ansatz}, measuring its energy on the quantum computer. Then, a new quantum state is prepared with a new set of parameters, selected by any classical optimization algorithm, and the process is repeated until the energy is minimized. If the energy is carried out perfectly, we approximate this value to the ground-state energy. \cite{VQEs}

Among the hybrid quantum-classical algorithms developed to address quantum simulation challenges, the \emph{Variational Quantum Eigensolver (VQE)} has gained particular relevance. That aims to solve for the ground state of a given Hamiltonian represented as a linear combination of Pauli terms, with an ansatz circuit where the number of parameters to optimize over is polynomial in the number of qubits. \cite{ibm_vqe_tutorial}

\subsection{Fundamental Principles and Stages of the Algorithm}
\label{subsec:vqe_principles_stages}

VQE is founded on the \textbf{variational principle} of quantum mechanics, which states that the expectation value of the energy 
\(\,E(\vec{\theta})\)
for any normalized trial state \(\,\ket{\psi(\vec{\theta})}\) is always an upper bound to the true ground-state energy \(E_0\):
\[
E(\vec{\theta}) 
= \bra{\psi(\vec{\theta})} \hat{H} \ket{\psi(\vec{\theta})}
\;\;\geq\; E_0.
\]

Because \(E(\vec{\theta})\) depends on a set of parameters \(\vec{\theta}\), the VQE iteratively updates these parameters to minimize the resultant energy. Once no further reduction in \(E(\vec{\theta})\) is possible, the algorithm identifies that we reached the best approximation of \(E_0\).
\begin{figure}[H]
    \centering
    \includegraphics[width=0.8\textwidth]{img/VQE_flowchart.png}
    \caption{VQE pipline\cite{VQE_review}.}
    \label{fig:vqe_flowchart}
\end{figure}

In this image, we find a flowchart of a VQE algorithm. Below, we describe the steps of the algoritm:


\begin{enumerate}
    \item \textbf{Hamiltonian Construction and representation:} 
    The first step of the algorithm is define the sytem for which we want to calculate the ground state. For that, we construc the Hamiltonian of the system. Constructing the Hamiltonian involves finding specific operators with their own weights.
    
    \item \textbf{Encoding of operators:}
    The next step, is to encode Hamiltonian operators, on quantum computers, we can only measure observables expressed in Pauli basis, due to the two-lavel nature of qubits. In first quantization, the operators can be directly translated into spin operators. In second quantization, the operators are expressed as a linear combination of fermionic operators.

    \item \textbf{Measurement of the Expected Energy:}
    The expectation value 
    \(\langle \psi(\vec{\theta}) \mid \hat{H} \mid \psi(\vec{\theta}) \rangle\)
    is obtained by measuring the appropriate Pauli operators on the quantum device, typically requiring multiple circuit executions due to non-commuting terms.

    \item \textbf{Ansatz and State Preparation:}
    Once we have prepared the Hamiltonian such that we can measure his expectation value on a quantum device, we can proceed to the preparation of the wavefunction. To achieve this, we need to decide the parametrized quantum circuit, denoted as ansatz. One of the main aspects that we have to take into account, is that the ansatz must be sufficiently expressive to guarantee that it can appropriately aproximate the ground state wavefunction.

    \item \textbf{Parameter Optimization:}
    Finally, the parameters of the ansatz used need to be updated iteratively until convergance. The measured energy serves as a cost function, we repeat the process until the convergence criteria is achieved.

    We can define the convergence criteria as:
    \[
    \left|\,E\bigl(\vec{\theta}_{k+1}\bigr) - E\bigl(\vec{\theta}_{k}\bigr)\right|
    \;<\;\delta,
    \]
    where \(\delta\) is a small threshold for energy differences. The final set \(\vec{\theta}^{*}\) yields an approximate ground-state wavefunction and energy.
\end{enumerate}

\subsection{Advantages and Challenges}
\label{subsec:vqe_challenges}

VQE has garnered extensive interest in fields like \textbf{quantum chemistry}, \textbf{materials science}, and even combinatorial optimization problems mapped to Hamiltonians. Its hybrid structure allows leveraging near-term quantum devices (with limited qubits and gate depths) while offloading resource-intensive tasks---such as parameter updates---to classical computers. 

Despite its promise, VQE faces several practical challenges:

\begin{itemize}
    \item \textbf{Noise and Decoherence:} 
    Real quantum devices suffer from errors that deteriorate the fidelity of prepared states and measurements, requiring error-mitigation strategies and noise-aware ansatz designs.

    \item \textbf{Barren Plateaus:}
    High-dimensional parameter spaces can contain large regions where gradients vanish, complicating the search for global minima.

    \item \textbf{Measurement Overhead:}
    Decomposing a Hamiltonian into many Pauli terms demands running multiple circuits, increasing sampling time and exposure to hardware noise.

    \item \textbf{Circuit Depth:}
    Accurate ans\"atze for complex systems may require deep circuits that quickly exceed the coherence times of current quantum processors.
\end{itemize}

\subsection{Outlook in Quantum Simulation}
\label{subsec:vqe_outlook}

Continuous improvements in hardware (qubit quality, gate fidelity, and error-correction schemes) and software (advanced ans\"atze, better optimizers, error mitigation) keep driving VQE toward practical applications. Recent strategies such as \emph{ADAPT-VQE}\cite{ADAPTVQE}, which determines a quasi-optimal ansatz with the minimal number of operators necessary to achieve the desired accuracy. The key idea is to systematically add fermionic operators one\-at\-a\-time, such that the maximal amount of correlation energy is captured at each step.

As the number of qubits grows and quantum hardware evolves, VQE is likely to become a central approach for tackling classically intractable problems in quantum chemistry, materials science, and beyond. Its flexible hybrid nature will continue to serve as a testbed for new optimization algorithms, ansatz designs, and measurement strategies, bridging current \emph{Noisy Intermediate-Scale Quantum (NISQ)} devices with the longer-term ambition of fault-tolerant quantum computing.

\section{Ans\"{a}tze}
In the context of variational quantum algorithms and quantum chemistry, an \textbf{ansatz} is a carefully chosen, often physically motivated, parametric form of the quantum state used to approximate the ground state of a system described by a given Hamiltonian. The term \textit{ansatz} originates from the German word \textquotedblleft approach\textquotedblright\ or \textquotedblleft initial guess,\textquotedblright\ and it reflects the central idea that we propose a functional form for the wavefunction and then we optimize the parameters in search of the lowest possible energy. 

Within the framework of the Variational Quantum Eigensolver (VQE), the ansatz is implemented as a parameterized quantum circuit whose gates depend on a set of continuous variables $\vec{\theta}$. It is used to produce the trial state, with which the Hamiltonian can be measured to obtain an estimate energy, $E(\vec{\theta})$, which is then iteratively minimized by a classical optimizer. The key aspects of the ansatz are the expressiveness and trainability of the chosen ansatz. The expressibility is the ability of the ansatz to span a large class of states, while the trainability describes the practical ability of the ansatz to be optimized using techniques tractable on quantum devices 
in the Hilbert space. 

Several ans\"{a}tze have been proposed to achieve a balance between accuracy and computational cost. Below, we summarize the most relevant approaches, highlighting their theoretical underpinnings and current usage in quantum simulation\cite{VQE_review}.

\subsection{Hartree--Fock-based Ans\"{a}tze (Classical Reference)}
A historically important \textquotedblleft classical\textquotedblright\ ansatz in quantum chemistry arises from the \textbf{Hartree--Fock (HF)}\cite{wiki:hartree_fock_method} approximation. In this method, the total wavefunction is assumed to be a single Slater determinant constructed from one-particle orbitals. Although it captures the fundamental antisymmetry required by the Pauli exclusion principle, it neglects most of the electron correlation. Post-HF methods, such as Configuration Interaction (CI), Many-Body Perturbation Theory (MPn), and Coupled Cluster (CC), then build on this reference state by introducing additional terms that account for electron correlation.\cite{wiki:post_hartree_fock}  

\begin{itemize}
    \item \textbf{Configuration Interaction (CI):} Expands the wavefunction in a basis of Slater determinants (excitations) beyond the HF reference. The Full CI approach is exact within the chosen basis but scales exponentially with system size. Truncated CI methods (CIS, CID, CISD, etc.) reduce the computational cost but still grow quickly with system size.

    \item \textbf{Coupled Cluster (CC):} CC is a post Hartree-Fock method that aims at recovering a portion of electron correlation energy by evolving an initial wave function (usually the Hartree-Fock wave function) under the action of parametrized excitation operators. Formally written as 
    \[
       |\Psi_{\mathrm{CC}}\rangle = e^{\hat{T}}|\Phi_{\mathrm{HF}}\rangle,
    \]
    where $\hat{T}$ is the sum of cluster excitation operators (singles, doubles, triples, etc.). Although Coupled Cluster with Singles and Doubles (CCSD) is often accurate, further inclusion of triples and higher excitations can be required for strongly correlated systems.
\end{itemize}

In the context of classical methods, these \textbf{ans\"{a}tze} serve as trial wavefunctions whose coefficients are optimized using high-performance classical algorithms. Their conceptual basis---constructing physically motivated trial states that capture crucial features of the system---carries over into quantum computing.

\subsection{Unitary Coupled Cluster (UCC)}
A key adaptation of the Coupled Cluster theory to quantum computing is the \textbf{Unitary Coupled Cluster (UCC)} ansatz. It modifies the standard CC exponential by making it explicitly unitary:

\[
|\Psi_{\mathrm{UCC}}\rangle 
= e^{\hat{T}(\vec{\theta}) - \hat{T}^\dagger(\vec{\theta})} \,|\Phi_{\mathrm{HF}}\rangle,
\]

where $\hat{T}(\vec{\theta})$ is typically truncated to include only single and double excitations (UCCSD). This approach guarantees that the resulting operator is unitary, which is crucial for hardware implementations in quantum computing since all gates must be unitary transformations. 

\begin{itemize}
    \item \textbf{UCCSD (Singles and Doubles):} The most widespread version of UCC is truncated at single and double excitations:
    \[
        \hat{T}(\vec{\theta}) = \sum_{i,a} \theta_{i}^{a} \hat{a}_a^\dagger \hat{a}_i 
          \,\,+\,\, 
          \sum_{i,j,a,b} \theta_{i,j}^{a,b} \hat{a}_a^\dagger \hat{a}_b^\dagger \hat{a}_j \hat{a}_i
          \,\,+\, \cdots
    \]
    Here, $i, j$ denote occupied orbitals and $a, b$ virtual (unoccupied) orbitals. By exponentiating both $\hat{T}$ and $\hat{T}^\dagger$, the wavefunction stays normalized. However, the circuit depth can become large since implementing the exponential of a sum of non-commuting operators requires a trotterization or related approximation.

    \item \textbf{ADAPT-VQE and Variants:} To mitigate the high circuit cost, variants like \textit{ADAPT-VQE} build up a UCC-type ansatz incrementally, selecting only those excitation operators that most significantly lower the energy at each step. This adaptive approach reduces the number of gates needed and often converges faster.
\end{itemize}

\section{Optimizers}
\label{sec:optimizers}

In quantum simulation algorithms such as the Variational Quantum Eigensolver (VQE), optimizers constitute a key element of the hybrid quantum-classical workflow. Their primary objective is to minimize the cost function
\[
E(\vec{\theta}) = \langle \psi(\vec{\theta}) \mid \hat{H} \mid \psi(\vec{\theta}) \rangle,
\]
where \(\hat{H}\) represents the Hamiltonian of the system under study and \(\ket{\psi(\vec{\theta})}\) is a parameterized quantum state, often referred to as the \emph{ansatz}. After each quantum measurement, the optimizer updates the parameter vector \(\vec{\theta}\) to guide the system toward the ground-state energy. This process is iterated until convergence, balancing the capabilities of quantum hardware with classical numerical techniques.

In the following subsections, we are going to describe, all the optimizers used in this project. While all these optimizers share the goal of efficiently navigating the parameter space, they differ in how they incorporate gradients, memory of past iterations, and adjustments of the learning rate.\cite{OptimizationAlgorithms} 

\subsection{Gradient Descent (GD)}
\label{subsec:gd}
\emph{Gradient Descent} is one of the most fundamental methods for continuous optimization. At each iteration, it updates parameters by moving them in the direction opposite to the gradient of the cost function:
\[
\vec{\theta}_{k+1} = \vec{\theta}_{k} 
- \eta \nabla_{\vec{\theta}} E(\vec{\theta}_{k}),
\]
where \(\eta\) is the learning rate and \(\nabla_{\vec{\theta}} E(\vec{\theta})\) denotes the gradient of the cost function with respect to the parameters. Gradient Descent can converge slowly or get trapped in local minima when dealing with complex or high-dimensional landscapes, making it less efficient if used alone in large-scale molecular simulations.

\subsection{Momentum Optimizer}
\label{subsec:momentum}
The \emph{Momentum} method extends standard Gradient Descent by incorporating a velocity term that accumulates a fraction of previous updates. This approach can mitigate oscillations and speed up convergence in regions with shallow gradients. The update rule is given by:
\[
\begin{aligned}
\vec{v}_{k+1} &= \gamma \vec{v}_{k} 
+ \eta \nabla_{\vec{\theta}} E(\vec{\theta}_{k}),\\
\vec{\theta}_{k+1} &= \vec{\theta}_{k} - \vec{v}_{k+1},
\end{aligned}
\]
where \(\gamma\) (typically between 0.9 and 0.99) is the momentum coefficient that controls how much past gradients influence the current update. This optimizer often accelerates learning in practice by effectively smoothing noisy or rapidly changing gradients.

\subsection{Nesterov Momentum Optimizer (NMomentum)}
\label{subsec:nmomentum}
\emph{Nesterov Momentum}, sometimes referred to as \emph{Nesterov’s Accelerated Gradient (NAG)}, refines the idea of Momentum by anticipating the next position of the parameters before computing the gradient. Concretely, one computes the gradient at \(\vec{\theta} + \gamma \vec{v}\) rather than at \(\vec{\theta}\) only. The updates become:
\[
\begin{aligned}
\vec{v}_{k+1} &= \gamma \vec{v}_{k} 
+ \eta \nabla_{\vec{\theta}}
E\!\big(\vec{\theta}_{k} + \gamma \vec{v}_{k}\big),\\
\vec{\theta}_{k+1} &= \vec{\theta}_{k} - \vec{v}_{k+1}.
\end{aligned}
\]
By \textquoteleft looking ahead\textquoteright\ in the direction of the velocity term, Nesterov Momentum tends to achieve smoother convergence and better performance on problems with numerous local minima or saddle points, which are common in complex quantum simulations.

\subsection{RMSProp}
\label{subsec:rmsprop}
\emph{RMSProp} is a gradient-based optimizer that adaptively tunes the learning rate for each parameter by normalizing the gradient through a moving average of its recent magnitudes. This helps address issues of vanishing or exploding gradients, which can be particularly troublesome in variational circuits of moderate or large depth. Its core update equations are:
\[
\begin{aligned}
E[\nabla^2_{\vec{\theta}}]_{k} &= \beta \, E[\nabla^2_{\vec{\theta}}]_{k-1}
+ (1-\beta)\,\nabla_{\vec{\theta}} E(\vec{\theta}_{k})^2,\\
\vec{\theta}_{k+1} &= \vec{\theta}_{k} 
- \eta \, \frac{\nabla_{\vec{\theta}} E(\vec{\theta}_{k})}
{\sqrt{E[\nabla^2_{\vec{\theta}}]_{k}} + \epsilon},
\end{aligned}
\]
where \(0 < \beta < 1\) is a decay factor controlling the smoothing effect, and \(\epsilon\) is a small constant ensuring numerical stability.

\subsection{Adagrad}
\label{subsec:adagrad}
\emph{Adagrad} is an early approach to adaptive learning rates, designed to handle sparse or highly nonuniform gradients. It individually scales the updates by the inverse square root of the cumulative sum of gradients:
\[
\vec{\theta}_{k+1} 
= \vec{\theta}_{k} 
- \frac{\eta}{\sqrt{\sum_{i=1}^{k} 
  \nabla_{\vec{\theta}} E(\vec{\theta}_{i})^2} + \epsilon}
  \,\nabla_{\vec{\theta}} E(\vec{\theta}_{k}).
\]
This mechanism allows parameters with small but consistent gradients to receive larger updates, which can be helpful in certain quantum chemistry models where specific Hamiltonian terms dominate.

\subsection{Adam}
\label{subsec:adam}
\emph{Adam (Adaptive Moment Estimation)} has emerged as one of the most widely used optimizers in machine learning and, increasingly, in quantum algorithms. It combines Momentum-like accumulations of the first moment of gradients (i.e., the mean) with an RMSProp-like treatment of the second moment (i.e., the uncentered variance). The update rules are:
\[
\begin{aligned}
\vec{m}_{k+1} &= \beta_1 \vec{m}_{k} 
+ (1-\beta_1) \nabla_{\vec{\theta}} E(\vec{\theta}_{k}), \\
\vec{v}_{k+1} &= \beta_2 \vec{v}_{k} 
+ (1-\beta_2) \nabla_{\vec{\theta}} E(\vec{\theta}_{k})^2, \\
\hat{\vec{m}}_{k+1} &= \frac{\vec{m}_{k+1}}{1-\beta_1^{k+1}}, \quad
\hat{\vec{v}}_{k+1} = \frac{\vec{v}_{k+1}}{1-\beta_2^{k+1}}, \\
\vec{\theta}_{k+1} &= \vec{\theta}_{k} 
- \eta \,\frac{\hat{\vec{m}}_{k+1}}
{\sqrt{\hat{\vec{v}}_{k+1}} + \epsilon},
\end{aligned}
\]
where \(0 < \beta_1, \beta_2 < 1\) are decay hyperparameters controlling how quickly the estimates of the first and second moments adjust. Adam’s blend of adaptive step sizes and momentum often yields robust performance, even when the cost landscape is noisy or irregular, as is typical in quantum simulations.

\subsection{Quantum Natural Gradient (QNG)}
\label{subsec:qng}
Unlike classical optimizers that rely on Euclidean metrics in parameter space, \emph{Quantum Natural Gradient (QNG)} specifically incorporates the \emph{Fubini--Study} metric, capturing how small changes in the parameters affect the underlying quantum state. This method is described in many \textit{Penny Lane} resources. We can also find the theoretical foundations of this optimizer.\cite{pennylane_qngoptimizer} Conceptually, the update rule can be written as:
\[
\vec{\theta}_{k+1}
= \vec{\theta}_{k}
- \eta \, \mathcal{F}^{-1} \nabla_{\vec{\theta}} E(\vec{\theta}_{k}),
\]
where \(\mathcal{F}\) represents the \emph{quantum Fisher information matrix}, a matrix encoding the local geometry of the parameterized state. Computing \(\mathcal{F}\) can be more demanding than classical gradients, but for many quantum chemistry, the improved efficiency justifies this added cost.

\subsection{Importance of Optimizers in Quantum Simulation}
\label{subsec:importance_optimizers}
Optimizers bridge the gap between quantum hardware and classical processing by iteratively refining the variational parameters to minimize the expectation value of the Hamiltonian. They must contend with challenges specific to quantum simulation, such as measurement noise, limited qubit counts, and complex cost landscapes characterized by local minima and barren plateaus. Properly choosing and tuning the optimizer is paramount to achieving accurate, resource-efficient simulations. By harnessing the distinctive advantages of adaptive and momentum-based methods---as well as more specialized quantum-aware techniques like QNG---one can significantly improve the speed and reliability of variational algorithms, thereby pushing the capabilities of quantum simulation closer to practical applications in molecular modeling.



% Methodology chapter. Please replace "methodology.tex" entirely with your own content written in the desired language.
%%%% PLEASE REPLACE ENTIRELY WITH YOUR OWN CONTENT %%%%


\chapter{Methodology / project development}

In this chapter, the methodology used in the completion of the work will be detailed. Its aim is to offer a thorough account of the approaches and techniques used, ensuring replicability and academic rigor. It will not only cover the research methods and measurement techniques employed but will also delve into the specifics of software and hardware development. Whether the project involves qualitative analysis, quantitative measurements, computational modeling, or physical prototyping, this chapter should elucidate how each component contributes to the overall objectives.

In addition to describing the methods themselves, the chapter will also provide justifications for why specific methods were chosen over others. For example, it may explain the choice of a particular programming language, statistical test, or experimental setup. The chapter will also address the limitations of the methodology and how these have been mitigated or accounted for. Readers should come away with a clear understanding of how the project's development has been carried out, why certain choices were made, and how these methods serve to fulfill the initially established objectives.

\section{Tools and Frameworks Selection} 
\subsection{Framework Selection}
To make the decision on which framework to use, we compared the documentation of the two quantum simulation frameworks available in the market: PennyLane and Qiskit. These are the most comprehensive frameworks with similar features available at the time of creating this project. After reviewing the documentation, we ultimately chose to use PennyLane for two reasons.

The first reason was the amount of documentation related to quantum simulation. Once we started looking into how others were using these resources, we realized that in the field of molecular simulation, the existing documentation—both theoretical and especially practical—was substantially greater. This provided us with more examples to begin developing our project.

The second reason for our choice was the frequent major changes implemented by Qiskit. We realized that while Qiskit is a tool that promises to be very good, it has historically undergone significant structural changes.

For these reasons, this project has been developed using the PennyLane framework. Below, we will observe how the project has been developed and explain the reasons behind the decisions made.

\section{Project Structuring}

After deciding on the interface to use and implementing the first version of the code, we decided to reorganize the project to make it more precise and modular. This structure offers the possibility to easily add more lines of code and functionalities.
\subsection{Code Organization}
\begin{ProjectStructure}
  \texttt{quantum\_simulation\_project/}
  \begin{itemize}[label={}, left=1em]
      \item \texttt{config/}
      \begin{itemize}[label={}, left=2em]
          \item \texttt{config\_functions.py}: Configuration functions for the project.
          \item \texttt{molecules.json}: Molecule data.
          \item \texttt{\_\_pycache\_\_}: Python cache files.
      \end{itemize}
      \item \texttt{main.py}: Main program file.
      \item \texttt{modules/}
      \begin{itemize}[label={}, left=2em]
          \item \texttt{ansatz\_preparer.py}: Quantum ansatz preparation.
          \item \texttt{hamiltonian\_builder.py}: Molecular Hamiltonian construction.
          \item \texttt{molecule\_manager.py}: Molecular data management.
          \item \texttt{opt\_mol.py}: Molecular optimization.
          \item \texttt{optimizer.py}: Optimization algorithms.
          \item \texttt{visualizer.py}: Visualization tools.
          \item \texttt{\_\_pycache\_\_}: Python cache files.
      \end{itemize}
      \item \texttt{temp\_results\_autograd/}
      \begin{itemize}[label={}, left=2em]
          \item \texttt{energy\_evolution.png}: Energy evolution graph.
          \item \texttt{filtered\_report\_autograd.txt}: Filtered results report.
          \item \texttt{final\_geometries\_3D.png}: Image of the final 3D geometries.
          \item \texttt{nuclear\_coordinates.png}: Nuclear coordinates.
          \item \texttt{output.txt}: Program data output.
          \item \texttt{profile\_output\_autograd.txt}: Autograd profile output.
      \end{itemize}
      \item \texttt{test/}: Directory for tests.
  \end{itemize}
\end{ProjectStructure}

This file serves as the entry point of the program. It initializes the initial conditions, sets up the molecule's configuration, defines optimization options, and orchestrates the execution of the quantum simulation and geometry optimization process. It is the first file to execute when running the application.

\textbf{Auxiliary Modules (\texttt{modules/})}:  
This directory contains various modules that encapsulate the core logic and operations of the project:
\begin{itemize}
    \item \texttt{ansatz\_preparer.py}: Includes functions for constructing the quantum circuit from the variational ansatz, preparing states, and applying excitations.
    \item \texttt{hamiltonian\_builder.py}: Responsible for generating the molecular Hamiltonian based on nuclear coordinates and the selected electronic basis.
    \item \texttt{molecule\_manager.py}: Manages information related to the molecule, such as its charge, multiplicity, number of electrons, and required orbitals.
    \item \texttt{opt\_mol.py}: Functions that orchestrate the molecular optimization process, invoking the optimizer and recording the simulation's progress.
    \item \texttt{optimizer.py}: Implements the optimization logic, combining routines for evaluating the cost (energy) with the selected optimization algorithms, updating parameters, and geometries.
    \item \texttt{visualizer.py}: Generates graphical outputs and reports that display the evolution of energy, nuclear coordinates, and other relevant data during optimization.
\end{itemize}

\textbf{Temporary Results Directories (\texttt{temp\_results\_autograd})}:  
    Several directories with names starting with \texttt{temp\_results\_autograd} store output data, time logs, and visualizations generated for different simulations and configurations. For documentation purposes, these directories are treated as a single repository for temporary results.

\textbf{Dependencies (\texttt{requirements.txt})}:  
    This file lists the required libraries and their versions to reproduce the project's execution environment. Keeping it updated ensures reproducibility of the simulation across different systems, facilitating the installation of necessary dependencies.


\subsection{Version Control}

Version control was managed using Git, allowing detailed tracking of changes and facilitating continuous collaboration with the supervisors on the project. Primarily, at the start of the project, a single branch was used to develop the project and explore the framework's possibilities. Once a stable version was achieved, branches were created to conduct tests and develop new functionalities. The first branches created were for the different interface versions. In each branch, the code was refined so that the same code would run across the various interfaces. Finally, only the interface changes that proved most suitable for the project were merged back into the main branch.

\section{Development and Implementation}

The \textit{Variational Quantum Eigensolver} (VQE) was chosen as the primary method to estimate the ground state energy of the studied quantum system. VQE combines limited quantum processing (measurements and applicability in moderately deep circuits) with classical optimization techniques. Its selection is justified by:

\begin{itemize}
    \item \textbf{Suitability for NISQ devices:} VQE is particularly well-suited for noisy intermediate-scale quantum (NISQ) devices, as it requires circuits of relatively low depth.
    \item \textbf{Flexible Ansatz:} It allows the use of various adaptive variational ansätze that capture essential electronic correlations.
    \item \textbf{Direct coupling to classical optimizers:} The VQE cost function (the expected energy) can be minimized with a wide range of classical methods, making it easy to experiment with different optimizers.
\end{itemize}

The core principle of VQE is the variational theorem, which guarantees that the expected energy of the ansatz is always an upper bound to the true ground state energy. By optimizing the ansatz parameters, the algorithm progressively approaches the actual energy minimum. We have already explained the concept of VQE in the state of the art chapter; now we will explain how we have implemented it in our project and how we have integrated it.

\paragraph{Principle of VQE:}

The VQE is based on the variational principle, which states that the expected energy of any approximate state \( |\psi(\theta)\rangle \) is always greater than or equal to the real ground state energy \( E_0 \):

\[
E(\theta) = \langle \psi(\theta) | H | \psi(\theta) \rangle \geq E_0
\]

We have already discussed this concept, but it is necessary to emphasize it as it is the foundation of the entire algorithm. The idea is to find the parameters \( \theta \) that minimize the expected energy, thereby approaching the real value of the ground state energy.

\bigskip

Next, we will detail how the VQE is implemented in our project, explaining how each component has been developed for trying to achieve the best performance and results.
\subsection{Hamiltonian Construction Process}

\begin{enumerate}
    \item \textbf{Definition of Molecular Geometry:}
    
    The geometry is specified by the atomic symbols and the Cartesian coordinates of each atom in the molecule:
    
    \begin{tcblisting}{colback=gray!5!white,colframe=gray!75!black,listing only,
      title=Geometry Definition, fonttitle=\bfseries, breakable, enhanced jigsaw, leftupper=8mm,
      listing options={language=Python, basicstyle=\ttfamily\small,
      showstringspaces=false, numbers=left, numberstyle=\footnotesize, stepnumber=1, numbersep=8pt}}
symbols = ['H', 'H']
x_init = np.array([0.0, 0.0, 0.0, 0.0, 0.0, 0.74]) 
    \end{tcblisting}
    
    
    \item \textbf{Hamiltonian Construction:}
    
    The \texttt{build\_hamiltonian} function generates the molecular Hamiltonian using PennyLane's functions:
    
    \begin{tcblisting}{colback=gray!5!white,colframe=gray!75!black,listing only,
      title=Hamiltonian Build, fonttitle=\bfseries, breakable, enhanced jigsaw, leftupper=8mm,
      listing options={language=Python, basicstyle=\ttfamily\small,
      showstringspaces=false, numbers=left, numberstyle=\footnotesize, stepnumber=1, numbersep=8pt}}
def build_hamiltonian(x, symbols, charge=0, mult=1, basis_name='sto-3g'):
    x = np.array(x)
    coordinates = x.reshape(-1, 3)
    hamiltonian, qubits = qml.qchem.molecular_hamiltonian(
        symbols, coordinates, charge=charge, mult=mult, basis=basis_name
    )
    h_coeffs, h_ops = hamiltonian.terms()
    h_coeffs = np.array(h_coeffs)
    hamiltonian = qml.Hamiltonian(h_coeffs, h_ops)
    return hamiltonian
    \end{tcblisting}
    \textbf{Note:}
    A basis set, such as 'sto-3g', is selected, which is a predefined set of basis functions to represent atomic orbitals in a simplified manner. This makes the simulation more efficient.
    
    \textbf{Function Description:}
    
    This function generates the qubit Hamiltonian of a molecule by transforming the electronic Hamiltonian in second quantization into the Pauli matrix framework. Additionally, it allows for the incorporation of net charge effects, spin multiplicity, and an active space defined by a specific number of electrons and orbitals, optimizing the quantum simulation of molecular systems.
    
\end{enumerate}

\subsection{Adaptive Ansatz Construction and Operator Selection}

As we said before, the adaptive ansatz construction builds upon the conventional variational approach by strategically selecting only those excitations that offer the most significant energy reductions. Instead of starting from a large, fixed set of parameters, the algorithm begins with the Hartree-Fock state and incrementally introduces new excitations based on their calculated impact on lowering the system’s energy. This methodology provides both theoretical and practical advantages in handling the complexity of the solution space.

The selection process begins with a predefined \textit{operator pool}, typically composed of single and double excitation operators relevant to the molecular system. At the start of the procedure, no variational parameters are assigned, and the system is initialized in the reference Hartree-Fock state. At each iteration, the algorithm evaluates the energy gradients associated with adding each operator from the pool:

\begin{enumerate}
    \item \textbf{Gradient Calculation:} For every candidate operator $\hat{O}_i$ in the pool, the partial derivative of the energy with respect to the parameter controlling $\hat{O}_i$ is computed. This step identifies how sensitive the energy is to introducing that particular excitation.
    \item \textbf{Operator Ranking and Filtering:} All candidate excitations are ranked according to the absolute value of their gradients. Operators that produce negligible energy changes are discarded, while those offering substantial decreases are selected for inclusion.
    \item \textbf{Incremental Ansatz Growth:} The selected operator(s) is then added to the ansatz. A new parameter is introduced and optimized, increasing the dimensionality of the parameter space \emph{only where it matters}. This targeted expansion ensures that each additional parameter contributes meaningfully to lowering the energy.
    \item \textbf{Pool Update and Iteration:} After adding the chosen operators, the process repeats. The operator pool is re-examined at subsequent steps, but it now excludes previously chosen operators unless they are included as parameterized parts of the ansatz. Over multiple iterations, the ansatz evolves adaptively, honing in on the most relevant subset of excitations.
\end{enumerate}

A simplified code snippet, consistent with the project’s structure, may appear as follows:

\begin{tcblisting}{colback=gray!5!white,colframe=gray!75!black,listing only,
  title=Adaptive Operator Selection, fonttitle=\bfseries, breakable, enhanced jigsaw, leftupper=8mm,
  listing options={language=Python, basicstyle=\ttfamily\small,
  showstringspaces=false, numbers=left, numberstyle=\footnotesize, stepnumber=1, numbersep=8pt}}
gradients = compute_operator_gradients(operator_pool, selected_excitations, params, hamiltonian, hf_state, dev, spin_orbitals)
selected_gate, max_grad_value = select_operator(gradients, operator_pool, convergence_threshold)
if selected_gate:
    selected_excitations.append(selected_gate)
    params = np.append(params, 0.0)  # Add new parameter for the chosen operator
    print(f"Added operator {selected_gate} with gradient {max_grad_value:.5e}")
else:
    print("No significant operators found. Convergence or local minimum reached.")
\end{tcblisting}

In this code, the \texttt{compute\_operator\_gradients} function evaluates each operator’s gradient, while the \texttt{select\_operator} function applies a filtering criterion based on a defined \texttt{convergence\_threshold}. Only the most promising excitation is incorporated into the ansatz at each step, ensuring a controlled and meaningful expansion of the parameter space.

In numerical experiments, this targeted approach has demonstrated:
\begin{itemize}
    \item \textbf{Faster Convergence:} Fewer parameters are introduced at each stage, allowing the optimizer to quickly reduce the energy without wading through irrelevant configurations.
    \item \textbf{Lower Resource Consumption:} By refining the search space, the quantum circuits remain relatively shallow, and classical optimization routines require fewer evaluations.
    \item \textbf{Scalability:} As molecular systems grow in complexity, the adaptive approach helps mitigate the exponential growth in parameter number, making it more feasible to handle larger systems within similar computational budgets.
\end{itemize}


\subsection{Cost Function Definition}
With the ansatz defined, the next step is to establish a cost function that evaluates the expected energy of the system given a set of parameters \(\theta\). In our implementation, this cost function is defined within \texttt{update\_parameters\_and\_coordinates} and calculates the expected value of the molecular Hamiltonian:
  
  \begin{tcblisting}{colback=gray!5!white,colframe=gray!75!black,listing only,
    title=Definition of the Cost Function, fonttitle=\bfseries, breakable, enhanced jigsaw, leftupper=8mm,
    listing options={language=Python, basicstyle=\ttfamily\small,
    showstringspaces=false, numbers=left, numberstyle=\footnotesize, stepnumber=1, numbersep=8pt}}
@qml.qnode(dev, interface=interface)
def cost_fn(params):
    prepare_ansatz(params, hf_state, selected_excitations, spin_orbitals)
    return qml.expval(hamiltonian)
  \end{tcblisting}
  
This function is essential for evaluating \(E(\theta)\). By calculating the expected value of the Hamiltonian, we can quantify how close our approximate state is to the true ground state.

\subsection{Mixed Electronic and Geometric Optimization Strategy}

In our project, we made a deliberate decision to integrate the refinement of both variational parameters \(\theta\) and nuclear coordinates \(\mathbf{X}\) into a single iterative loop. Rather than treating these aspects separately, as is common in sequential optimization approaches, we chose to couple the electronic and geometric optimizations from the start. By doing so, we ensured that each electronic update immediately reflects the current molecular geometry, and each geometric adjustment leverages the most recent and accurate electronic wavefunction. This approach was guided by our goal of finding the global minimum of the molecular system’s total energy more efficiently, ultimately steering the molecule towards its equilibrium geometry and ground state energy in a more coherent and resource-effective manner.

\paragraph{Rationale for Our Mixed Approach}
In many standard methodologies, the electronic structure is optimized at a fixed geometry before the geometry itself is optimized under that electronic configuration. We felt this separation was inefficient for our objectives, as it often leads to unnecessary iterations and suboptimal intermediate solutions. Recognizing that the electronic configuration and nuclear arrangement are intrinsically linked, we decided to continually update both in tandem. This choice was grounded in our understanding that capturing the interplay between the electronic distribution and molecular geometry as it unfolds would reduce computational overhead and converge more rapidly to a stable, physically meaningful equilibrium.

\paragraph{Our Iterative Optimization Steps}
\begin{enumerate}
    \item \textbf{Initialization:}
    We begin by defining the initial molecular geometry, nuclear coordinates \(\mathbf{X}_0\), and variational parameters \(\theta_0\). For instance:
    \begin{tcblisting}{colback=gray!5!white,colframe=gray!75!black,listing only,
      title=Initialization Example, fonttitle=\bfseries, breakable, enhanced jigsaw, leftupper=8mm,
      listing options={language=Python, basicstyle=\ttfamily\small,
      showstringspaces=false, numbers=left, numberstyle=\footnotesize, stepnumber=1, numbersep=8pt}}
symbols = ['H', 'H']
x_init = np.array([0.0, 0.0, 0.0, 0.0, 0.0, 0.74])
params = np.array([], requires_grad=True)
    \end{tcblisting}
    We selected these initial values based on known molecular structures and our desire to start from a configuration that is simple yet representative.

    \item \textbf{Molecular Hamiltonian Update:}
    At each iteration, we recalculate the molecular Hamiltonian \(\hat{H}(\mathbf{X})\) using the updated nuclear coordinates. This ensures that our cost function (the expected energy) accurately reflects the current geometry at every step, allowing us to maintain a closer link between the electronic solution and the evolving molecular shape.

    \item \textbf{Operator Gradient Calculation and Selection:}
    For the electronic part, we compute energy gradients with respect to a pool of candidate excitation operators. We then select and add only the operator with the highest gradient contribution to the ansatz. This choice was made to keep the parameter space as manageable as possible, focusing our resources on the excitations that yield the greatest energy reductions. By gradually expanding the ansatz in this manner, we prevent an explosion in complexity and retain a more controlled optimization path.

    \item \textbf{Nuclear Gradient Computation and Coordinate Updates:}
    To update the nuclear coordinates \(\mathbf{X}\), we decided to compute the nuclear gradients numerically using finite differences. Concretely, for each coordinate direction, we slightly perturb \(\mathbf{X}\) and measure the resulting change in energy:
    \begin{equation}
    \nabla_{\mathbf{X}} E(\theta, \mathbf{X}) \approx \frac{E(\theta, \mathbf{X} + \delta \mathbf{u}) - E(\theta, \mathbf{X} - \delta \mathbf{u})}{2\delta}.
    \end{equation}
    We adopted this finite-difference method to maintain flexibility and simplicity, and because it integrates seamlessly with our adaptive approach. After obtaining the gradients, we update the coordinates as:
    \begin{equation}
    \mathbf{X}_{\text{new}} = \mathbf{X}_{\text{old}} - \alpha \nabla_{\mathbf{X}}E(\theta, \mathbf{X}),
    \end{equation}
    where \(\alpha\) is a learning rate we chose based on preliminary tests to balance stability and speed of convergence. This direct control allowed us to refine the geometry efficiently without overshooting or requiring overly complex gradient estimations.

    \item \textbf{Geometric Convergence Validation:}
    We established convergence criteria tailored to our desired precision and computational budget. For energy convergence, we require that differences between subsequent iterations fall below a small threshold (e.g., \(10^{-8}\) Ha). For geometric convergence, we set clear tolerances on changes in atomic positions and interatomic distances. Meeting both criteria ensures the geometry is near equilibrium and that no further improvements are necessary. This choice reflects our priority for a stable and physically meaningful final configuration rather than just a numerically optimized one.

    \item \textbf{Visualization and Termination:}
    Throughout the process, we record the evolving energy, coordinates, and other metrics. We decided to include immediate visual feedback (such as energy vs. iteration plots and 3D molecular views) to quickly assess the effectiveness of our approach. This real-time monitoring helps us recognize potential issues early, adjust our strategies, and confirm when we have reached the desired convergence criteria. Once these conditions are met, we terminate the process, confident that the final geometry and electronic state are both well-optimized.
\end{enumerate}

\paragraph{Why Our Iterative Coupling Is More Efficient}
By integrating electronic and geometric optimizations, we reduce the redundant recalculations and needless complexity that can arise in standard sequential methods. In our approach, each update to the electronic parameters leverages a geometry that is already moving towards equilibrium, and each geometric refinement uses an increasingly accurate electronic state. This synergy not only shortens the pathway to convergence but also helps us avoid getting stuck in suboptimal configurations. In our experience, this mixed strategy has proven more efficient and more closely aligned with the underlying physics of the system, delivering faster, more stable, and more accurate results than what we would expect from a traditional, decoupled optimization approach.

\section{Limitations and Mitigation Measures}

Among the limitations of this approach are:
\begin{itemize}
    \item \textbf{Scalability:} As the system grows in the number of electrons and orbitals, the complexity of Hamiltonian construction and the excitation space increases exponentially.
    \item \textbf{Quantum Noise and Errors:} On real devices, noise affects measurement fidelity. Our work, primarily simulation-oriented, plans to integrate mitigation techniques in future studies.
    \item \textbf{Ansatz Choice:} Although the adaptive ansatz helps, there is no guarantee that the excitation selection is optimal. Future work might explore more complex heuristics.
\end{itemize}

To mitigate these issues, we opted for reduced basis sets, strategies such as re-initializing the optimizer when increasing the parameter space, and verifying convergence through multiple criteria (energetic and geometric).



%%%% RESULTS %%%%
% Results chapter. Please replace "results.tex" entirely with your own content written in the desired language.
\chapter{Results}
This chapter should encompass your data analysis and findings. Additionally, include relevant tables, figures, and citations to support your results and interpretations. Here is a suggested list of topics to discuss:
\section{Interface Comparison}
During the development of this project, various experiments were conducted to optimize the performance of a quantum molecular simulator. In this context, the choice of the interface for calculating cost functions is crucial, as these functions must be optimized to obtain the optimal state and geometry of each molecule. Typically, the \textit{JAX} interface is the most widely used due to its GPU acceleration capabilities, which generally provide greater efficiency and speed in searching for the minimum of the cost function.

However, when comparing two identical implementations—varying only the calculation interface—we observed a surprising result: the \textit{JAX}-based interface, which was executed on a GPU, was slower than the \textit{autograd}-based interface, which ran on a CPU. To confirm that this was not a coding error, we repeated the experiments on different machines, obtaining the same result. For this reason, we decided to continue with the \textit{autograd} interface due to its greater speed in our specific implementation.

\subsection{Optimization and Timing Logging}
To verify the performance differences between both interfaces, the same optimization was executed while varying only the interface. During these simulations, the execution times of the different functions in the code were recorded, with their values shown in Table~\ref{tab:comparison_gd}.

\begin{table}[H]
  \centering
  \scriptsize
  \resizebox{\textwidth}{!}{%
  \begin{tabular}{lcccccc}
  \toprule
  \textbf{Function} & 
  \textbf{[\texttt{Li}, \texttt{H}] \_autograd\_GD} & 
  \textbf{[\texttt{Li}, \texttt{H}] \_jax\_GD} & 
  \textbf{[\texttt{O}, \texttt{H}, \texttt{H}] \_autograd\_GD} & 
  \textbf{[\texttt{O}, \texttt{H}, \texttt{H}] \_jax\_GD} & 
  \textbf{[\texttt{H}, \texttt{H}] \_autograd\_GD} & 
  \textbf{[\texttt{H}, \texttt{H}] \_jax\_GD} \\
  \midrule
  Iteration 1  & 1246.7627 & 6251.3362 & 6942.21696 & 39003.89009 & 14.6908 & 40.5485 \\
  Iteration 2  & 1244.0321 & 6284.9812 & 6935.09353 & 40652.7751  & 14.7824 & 36.2979 \\
  Iteration 3  & 1247.4399 & 6348.2470 & 6920.68571 & 40386.29868 & 14.9931 & 38.6236 \\
  Iteration 4  & 1245.7256 & 6351.6437 & 6942.60873 & 40128.31419 & 15.0694 & 39.1362 \\
  Iteration 5  & 1244.8591 & 6349.5099 & 6965.70663 & 40657.74473 & 15.1512 & 40.6186 \\
  \midrule
  \textbf{Total Time} & 6228.8195 & 31585.7475 & 34706.3116 & 200829.0866 & 74.6870 & 195.2551 \\
  \textbf{build\_hamiltonian} & 29.1349   & 30.3407  & 78.4744   & 85.7532   & 0.4934   & 0.5329  \\
  \textbf{compute\_operator\_gradients} & 1533.0319 & 17497.4381 & 12563.3726 & 113792.7167 & 0.6313   & 11.9006 \\
  \textbf{update\_parameters\_and\_coordinates} & 4666.6385 & 14056.6795 & 22064.4153 & 86946.6021 & 73.5594 & 182.5271 \\
  \bottomrule
  \end{tabular}%
  }
  \caption{Execution times for different molecules and interfaces using \textit{Gradient Descent}.}
  \label{tab:comparison_gd}
\end{table}

Figure~\ref{fig:time_iterations} shows how the execution time evolves over the iterations for each molecule and type of interface. It is evident that, in all cases, the \textit{autograd}-based interface offers lower computation times than the \textit{JAX}-based interface.

\begin{figure}[H]
  \centering
  \includegraphics[width=0.8\textwidth]{img/time_iterations.png}
  \caption{Execution time of different molecules per iteration in the simulation.}
  \label{fig:time_iterations}
\end{figure}

\subsubsection{Comparative Performance Analysis}
To evaluate the performance of both interfaces in more detail, a table was created showing the percentage increase in execution time of \textit{JAX} compared to \textit{autograd}, as seen in Table~\ref{tab:jax_vs_auto}.

\begin{table}[H]
  \centering
  \scriptsize
  \resizebox{\textwidth}{!}{%
  \begin{tabular}{lcccc}
  \toprule
  \textbf{Metric} & 
  \textbf{LiH} & 
  \textbf{H\textsubscript{2}} & 
  \textbf{H\textsubscript{2}O} & 
  \textbf{Mean} \\
  \midrule
  \textbf{Total Time} & 80.28\% & 82.72\% & 61.75\% & 74.92\% \\
  \textbf{build\_hamiltonian} & 3.97\% & 8.49\% & 7.42\% & 6.63\% \\
  \textbf{compute\_operator\_gradients} & 91.24\% & 88.96\% & 94.70\% & 91.63\% \\
  \textbf{update\_parameters\_and\_coordinates} & 66.80\% & 74.62\% & 59.70\% & 67.04\% \\
  \bottomrule
  \end{tabular}%
  }
  \caption{Comparison of execution times between JAX and autograd interfaces for different molecules.}
  \label{tab:jax_vs_auto}
\end{table}

On average, it was observed that the \textit{JAX} interface exhibits a 74.92\% higher total execution time than the \textit{autograd} interface. Additionally, it is interesting to note that as the complexity of the molecule increases (with a greater number of atoms), the penalty percentage of \textit{JAX} tends to decrease. This behavior suggests that, in larger-scale problems, GPU acceleration could become more competitive, although it does not manage to outperform \textit{autograd} in this implementation.

\subsection{Computation Time per Function}
For a higher level of detail, the computation time was also measured for each part of the code where the interface change is introduced. Figure~\ref{fig:time_functions} shows the accumulated execution times in the main stages of the algorithm.

\begin{figure}[H]
  \centering
  \includegraphics[width=0.8\textwidth]{img/time_functions.png}
  \caption{Execution time in different parts of the code.}
  \label{fig:time_functions}
\end{figure}

The same pattern is maintained in all functions: the \textit{JAX} interface records higher execution times than \textit{autograd}. In particular, the gradient computation part \texttt{compute\_operator\_gradients} increases its execution time by 91.68\% when using \textit{JAX}. This difference is largely attributed to the overhead caused by data transfer between the CPU and GPU, especially when \textit{PennyLane} does not fully support generating the molecule's Hamiltonian directly on the GPU.

On the other hand, the function \texttt{update\_parameters\_and\_coordinates}, responsible for performing the molecular geometry and parameter optimization step, also shows a 67.04\% increase when using \textit{JAX}. Nevertheless, it is worth highlighting that as the problem grows in complexity, the \textit{JAX} interface gains some relative efficiency in gradient calculation; however, the time saved through GPU acceleration is offset by the continuous data transfer between CPU and GPU throughout the iterations.

\subsection{Conclusions}
Based on the obtained results, the \textit{autograd} interface demonstrated superior performance in terms of speed for our implementation of the quantum molecular simulator. Although the GPU is usually advantageous in larger-scale problems, the data transfer overhead and the lack of full support for Hamiltonian generation within the GPU reduced the efficiency of \textit{JAX}. For this reason, we ultimately chose to use the \textit{autograd} interface for the final implementation of the quantum molecular simulator.

\section{Ansatz Comparison}
One of the most important modifications that has significantly affected our code and its functionality has been the choice of the Ansatz. Our proposal has been to implement UCCSD, an Ansatz typically used for this type of simulation, as it enhances the simulation performance by achieving higher efficiency. The efficacy of this type of Ansatz has already been demonstrated, showing how it can improve simulation performance by producing more optimal quantum circuits without the need to create a specific Ansatz for the molecule being simulated. It is a fact that for each optimizer configuration and for each different molecule simulation, there exists an optimal quantum circuit that achieves the best performance. However, since our objective is to develop a program that can simulate various molecules with maximum performance, we observe that the best option is UCCSD. To illustrate how our simulation performance is improved, we have generated Ansätze with different levels of depth and compared their performance with that of a UCCSD Ansatz. We have simulated various configurations of classical Ansätze, and in all cases, the UCCSD Ansatz has achieved better performance. More complex and molecule-specific Ansätze could be tested, but it is unlikely that another type of Ansatz would outperform UCCSD in terms of performance improvement.

Below is a simulation with different Ansatz depths and varying numbers of iterations obtained directly from the simulation.

\begin{figure}[H]
  \centering
  \includegraphics[width=0.8\textwidth]{data/Anzatz/results_ansatz_lyers_dif_iterations/energy_vs_time_log_offset.png}
  \caption{Energy vs. Time for different Ansätze and iterations.}
  \label{fig:ansatz_layers_iterations}
\end{figure}

It is observed that the UCCSD Ansatz achieves the best performance in all simulations, regardless of the number of optimizations performed for each iteration. For greater clarity of the results, a simulation was conducted with only a single number of optimizations per iteration, and the performance of the different Ansätze was compared, providing a clearer view of how UCCSD achieves the best performance.

\begin{figure}[H]
  \centering
  \includegraphics[width=0.8\textwidth]{data/Anzatz/results_ansatz_lyers/energy_vs_time_log_offset.png}
  \caption{Energy vs. Time for different Ansätze.}
  \label{fig:ansatz_layers}
\end{figure}

Finally, to corroborate that the UCCSD Ansatz achieves the best performance, a simulation was conducted with a different molecule, and the performance of the various Ansätze was compared. In the following image, it can be seen that the UCCSD Ansatz achieves the best performance in all simulations.

\section{Optimizer}

Once the Ansatz was selected, the next step was to choose and configure the optimizer. For this purpose, preliminary simulations were conducted to identify the best configuration.

\subsection{Optimizer Selection}
To determine the most effective optimizer for different molecules, a set of simulations was designed to analyze the evolution of energy as a function of iterations. To achieve this, we developed a procedure that allowed executing the same molecules with various optimizers, each covering a range of \emph{step size} values. This approach enabled the selection of the most appropriate \emph{step size} for each case.

In the first phase, which executed 42 processes in parallel, the \emph{step size} offering the best performance was observed for each optimizer and molecule. Subsequently, the simulation was repeated using only those optimal \emph{step sizes}. The results are shown below.

\begin{figure}[H]
  \centering
  % First image
  \begin{subfigure}{0.32\textwidth}
    \includegraphics[width=\textwidth]{data/Optimizadores/final_results_H2/energy_evolution_log_offset.png}
    \caption{H2 simulation.}
    \label{fig:subimage1}
  \end{subfigure}
  % Second image
  \begin{subfigure}{0.32\textwidth}
    \includegraphics[width=\textwidth]{data/Optimizadores/final_results_LiH/energy_evolution_log_offset.png}
    \caption{LiH simulation.}
    \label{fig:subimage2}
  \end{subfigure}
  % Third image
  \begin{subfigure}{0.32\textwidth}
    \includegraphics[width=\textwidth]{data/Optimizadores/final_results_H20/energy_evolution_log_offset.png}
    \caption{H2O simulation.}
    \label{fig:subimage3}
  \end{subfigure}
  \caption{Optimal \emph{step size} selection for each optimizer and molecule.}
  \label{fig:three_images}
\end{figure}

To delve deeper into the performance of the different optimizers, the following tables present the final energy, total optimization time, and the number of iterations required to converge for each molecule.

\subsubsection{Molecule \(\mathrm{H_2O}\)}

\begin{table}[H]
\centering
\caption{Final energy and optimization time for \(\mathrm{H_2O}\) (without using FCI).}
\begin{tabular}{lccc}
\toprule
\textbf{Optimizer} & \textbf{Final Energy (Ha)} & \textbf{Total Time (s)} & \textbf{Iterations} \\
\midrule
Adam (0.5)       & -73.75990609 & 28454.49 & 10 \\
Adagrad (0.6)    & -73.93046296 & 28891.08 & 10 \\
NMomentum (0.5)  & -73.22152796 & 7497.83  &  1 \\
Momentum (0.2)   & \textbf{-74.03997489} & 68420.47 & 10 \\
RMSProp (0.5)    & -73.28585876 & 65470.02 & 10 \\
GD (0.5)         & -73.22152796 & \textbf{6486.79}  &  1 \\
QNG (0.5)        & -73.13971041 & 70716.22 & 10 \\
\bottomrule
\end{tabular}
\end{table}

In this molecule, the most \emph{negative} energy is provided by \textbf{Momentum (0.2)}, with \(-74.03997489\)\,Ha, but at the cost of a very high time (68420.47\,s). On the other hand, \textbf{GD (0.5)} requires only 6486.79\,s, although its energy value \((-73.22152796)\,\mathrm{Ha}\) is somewhat less negative.

\subsubsection{Molecule \(\mathrm{H_2}\)}

\begin{table}[h!]
\centering
\caption{Final energy and optimization time for \(\mathrm{H_2}\).}
\begin{tabular}{lccc}
\toprule
\textbf{Optimizer} & \textbf{Final Energy (Ha)} & \textbf{Total Time (s)} & \textbf{Iterations} \\
\midrule
Adam (0.1)       & -1.07655292  & 173.30  & 10 \\
Adagrad (0.25)   & -1.13469066  & 10.19   &  1 \\
NMomentum (0.25) & -1.13730600  & 62.84   &  4 \\
Momentum (0.3)   & \textbf{-1.13730605} & 100.43  &  6 \\
RMSProp (0.1)    & -1.13675411  & 33.14   &  2 \\
GD (0.3)         & \textbf{-1.13730605} & 41.77   &  3 \\
QNG (0.1)        & -1.13469066  & 18.40   &  1 \\
\bottomrule
\end{tabular}
\end{table}

In this case, the most negative values (\(-1.13730605\)\,Ha) correspond to both \textbf{Momentum (0.3)} and \textbf{GD (0.3)}. If speed is a priority, \textbf{GD (0.3)} only requires 41.77\,s, making it a very competitive option.

\subsubsection{Molecule \(\mathrm{LiH}\)}

\begin{table}[H]
\centering
\caption{Final energy and optimization time for \(\mathrm{LiH}\).}
\begin{tabular}{lccc}
\toprule
\textbf{Optimizer} & \textbf{Final Energy (Ha)} & \textbf{Total Time (s)} & \textbf{Iterations} \\
\midrule
Adam (0.1)       & -7.87024707 & 13444.57 & 10 \\
Adagrad (0.2)    & -7.80548501 & 977.86   &  1 \\
NMomentum (0.05) & \textbf{-7.87085783} & 13442.59 & 10 \\
Momentum (0.1)   & -7.75267240 & 13566.98 & 10 \\
RMSProp (0.15)   & -7.67650000 & 13671.22 & 10 \\
GD (0.02)        & -7.86422129 & 13449.49 & 10 \\
QNG (0.01)       & -7.85120449 & 13776.00 & 10 \\
\bottomrule
\end{tabular}
\end{table}

Here, the lowest value \((-7.87085783)\,\mathrm{Ha}\) is achieved with \textbf{NMomentum (0.05)}. Both NMomentum (0.05) and Adam (0.1) surpass \(-7.870\)\,Ha with very similar times (both around 13400\,s). Meanwhile, \textbf{Momentum (0.1)} proves to be slower and less effective in terms of energy.

In summary, \emph{there is no universally optimal optimizer}. The choice depends on the molecule, acceptable computational cost, and the goal in terms of final energy. \textbf{Momentum} usually reaches the lowest values, although sometimes with high time requirements. \textbf{GD}, on the other hand, converges quickly and offers results close to the most negative minima in several tests. After reviewing the obtained results, we decided to choose the Momentum optimizer, as it achieves the best convergence without requiring excessive computation time.

\subsection{Step Size Selection}
To determine the most suitable \textit{step size} range, we started with the optimal value identified during optimizer selection. Subsequently, additional simulations were designed with values close to the initial optimum to refine this selection. This section presents the evolution of energy as a function of iterations, considering different \textit{step size} values for each molecule.

\begin{figure}[H]
  \centering
  % First image
  \begin{subfigure}{0.32\textwidth}
    \includegraphics[width=\textwidth]{data/StepSize/results_H2/energy_evolution_log_offset.png}
    \caption{H$_2$ simulation.}
    \label{fig:step_size_h2}
  \end{subfigure}
  % Second image
  \begin{subfigure}{0.32\textwidth}
    \includegraphics[width=\textwidth]{data/StepSize/results_LiH/energy_evolution_log_offset.png}
    \caption{LiH simulation.}
    \label{fig:step_size_lih}
  \end{subfigure}
  % Third image
  \begin{subfigure}{0.32\textwidth}
    \includegraphics[width=\textwidth]{data/StepSize/results_H2O/energy_evolution_log_offset.png}
    \caption{H$_2$O simulation.}
    \label{fig:step_size_h2o}
  \end{subfigure}
  \caption{Energy evolution as a function of \textit{step size} for different molecules.}
  \label{fig:step_size_results}
\end{figure}

Below, the final tables of energy, total optimization time, and the number of iterations for each molecule are shown. Note that in this case, each row corresponds to a \textit{step size}, keeping the same optimizer (in our case, \textbf{Momentum}).

\subsubsection{Molecule \(\mathrm{H_2}\)}

\begin{table}[H]
\centering
\caption{Final energy and optimization time for \(\mathrm{H_2}\) varying the \textit{step size} (optimizer: Momentum).}
\begin{tabular}{lccc}
\toprule
\textbf{Step Size} & \textbf{Final Energy (Ha)} & \textbf{Total Time (s)} & \textbf{Iterations} \\
\midrule
0.10 & \(\mathbf{-1.13730605}\) & 113.23 & 10 \\
0.12 & \(\mathbf{-1.13730605}\) & \(\mathbf{92.55}\) & 10 \\
0.14 & \(\mathbf{-1.13730605}\) & 131.52 & 10 \\
0.16 & \(\mathbf{-1.13730605}\) & 101.46 & 10 \\
0.18 & \(\mathbf{-1.13730605}\) & 149.96 & 10 \\
0.20 & -1.13730604 & 128.31 & 10 \\
0.22 & \(\mathbf{-1.13730605}\) & 111.54 & 10 \\
0.24 & \(\mathbf{-1.13730605}\) & 123.74 & 10 \\
0.26 & \(\mathbf{-1.13730605}\) & 115.02 & 10 \\
\bottomrule
\end{tabular}
\end{table}

In \(\mathrm{H_2}\), almost all \textit{step sizes} reach the same most negative energy \(\bigl(-1.13730605\,\mathrm{Ha}\bigr)\). If we prioritize speed, the \textit{step size} \(0.12\) is the most convenient, as it achieves this value in only \(92.55\,\mathrm{s}\).

\subsubsection{Molecule \(\mathrm{H_2O}\)}

\begin{table}[H]
\centering
\caption{Final energy and optimization time for \(\mathrm{H_2O}\) varying the \textit{step size} (optimizer: Momentum).}
\begin{tabular}{lccc}
\toprule
\textbf{Step Size} & \textbf{Final Energy (Ha)} & \textbf{Total Time (s)} & \textbf{Iterations} \\
\midrule
0.10 & \(\mathbf{-74.68985518}\) & 48577.15 & 10 \\
0.13 & -74.46832273 & 74175.94 & 10 \\
0.15 & -74.25355192 & 74211.21 & 10 \\
0.18 & -74.67526296 & 75127.97 & 10 \\
0.20 & -74.03997489 & 73130.72 & 10 \\
0.24 & -73.99558215 & 56465.58 & 10 \\
0.26 & -74.68963596 & \(\mathbf{44970.94}\) & 10 \\
0.28 & -74.59437540 & 75555.08 & 10 \\
0.30 & -74.65314486 & 73575.44 & 10 \\
\bottomrule
\end{tabular}
\end{table}

For \(\mathrm{H_2O}\), the \textit{step size} \(0.10\) offers the most negative energy \(\bigl(-74.68985518\,\mathrm{Ha}\bigr)\), while \(0.26\) notably reduces the total time to approximately \(44970.94\,\mathrm{s}\), with a very competitive energy value \(\bigl(-74.68963596\,\mathrm{Ha}\bigr)\).

\subsubsection{Molecule \(\mathrm{LiH}\)}

\begin{table}[H]
\centering
\caption{Final energy and optimization time for \(\mathrm{LiH}\) varying the \textit{step size} (optimizer: Momentum).}
\begin{tabular}{lccc}
\toprule
\textbf{Step Size} & \textbf{Final Energy (Ha)} & \textbf{Total Time (s)} & \textbf{Iterations} \\
\midrule
0.02 & -7.86678430 & \(\mathbf{13915.43}\) & 10 \\
0.04 & -7.87699744 & 14001.06 & 10 \\
0.06 & -7.88014045 & 13916.11 & 10 \\
0.08 & -7.87944832 & 14144.26 & 10 \\
0.10 & -7.75267240 & 14098.79 & 10 \\
0.12 & -7.87992395 & 13937.66 & 10 \\
0.14 & -7.78203292 & 14063.16 & 10 \\
0.16 & \(\mathbf{-7.88118152}\) & 13931.81 & 10 \\
0.18 & -7.88116774 & 13998.38 & 10 \\
\bottomrule
\end{tabular}
\end{table}

In \(\mathrm{LiH}\), the \textit{step size} \(0.16\) achieves the most negative final energy \(\bigl(-7.88118152\,\mathrm{Ha}\bigr)\). If the goal is to minimize computation time, the \textit{step size} \(0.02\) stands out with only \(13915.43\,\mathrm{s}\).

Thus, it is concluded that the choice of \textit{step size} depends both on the energy sought to be optimized and the acceptable computation time. For \(\mathrm{H_2}\), for example, a \textit{step size} of \(0.12\) achieves an excellent compromise between final energy and speed; whereas, for \(\mathrm{H_2O}\) and \(\mathrm{LiH}\), one might opt for values that minimize energy (case \(\mathrm{H_2O}\) with \(0.10\) and \(\mathrm{LiH}\) with \(0.16\)) or for shorter convergence times (for example, \(0.26\) in \(\mathrm{H_2O}\) and \(0.02\) in \(\mathrm{LiH}\)).

\section{Numero de Iteraciones}

Finalmente, una de los procesos donde mas tiempo se consume en la simulación es en la actualización de parametros y de las coordenadas, en esta ultima prueba lo que queremos conseguir es poder optimizar el numero de iteraciones que se actualizan los parametros y las coordenadas antes de volver a calcular el Hamiltoniano. Para ello, hemos generado una serie de simulaciones con distintos numeros de iteraciones, y hemos observado la evolución de la energia en función de las iteraciones. La configuarción de los optimizadores, han sido con los valores optimos que hemos obtenido despues de haber hecho todas las pruebas, con el optimizador Momentum y los valores de step size optimos.

\begin{figure}[H]
  \centering
  % First image
  \begin{subfigure}{0.45\textwidth}
    \includegraphics[width=\textwidth]{data/NumIterations/results_H2/energy_evolution_log_offset.png}
    \caption{H$_2$ simulation.}
    \label{fig:num_iterations_h2}
  \end{subfigure}
  % Second image
  \begin{subfigure}{0.45\textwidth}
    \includegraphics[width=\textwidth]{data/NumIterations/results_LiH/energy_evolution_log_offset.png}
    \caption{LiH simulation.}
    \label{fig:num_iterations_lih}
  \end{subfigure}
  % Third image
  \begin{subfigure}{0.45\textwidth}
    \includegraphics[width=\textwidth]{data/StepSize/results_H2O/energy_evolution_log_offset.png}
    \caption{H$_2$O simulation.}
    \label{fig:num_iterations_h2o}
  \end{subfigure}
  \caption{Energy evolution as a function of \textit{number of iterations} for different molecules.}
  \label{fig:num_iterations_results}
\end{figure}

Igual que en los apartados anteriores, a continuación mostramos las tablas con los valores de energia final, tiempo total de optimización y el numero de iteraciones necesarias para converger para cada molécula, para poder asi analizar los resultados obtenidos.

\subsubsection{Molecule \(\mathrm{H_2}\)}

\begin{table}[H]
\centering
\caption{Final energy and optimization time for \(\mathrm{H_2}\) using Momentum optimizer with varying steps.}
\begin{tabular}{lccc}
\toprule
\textbf{Steps} & \textbf{Final Energy (Ha)} & \textbf{Total Time (s)} & \textbf{Difference from FCI (Ha)} \\
\midrule
5 & \(\mathbf{-1.13730605}\) & 101.86 & -0.019800 \\ 
10 & \(\mathbf{-1.13730605}\) & \(\mathbf{99.08}\) & -0.019800 \\ 
15 & \(\mathbf{-1.13730605}\) & 163.25 & -0.019800 \\ 
20 & \(\mathbf{-1.13730605}\) & 175.06 & -0.019800 \\ 
25 & \(\mathbf{-1.13730605}\) & 124.58 & -0.019800 \\ 
30 & \(\mathbf{-1.13730605}\) & 147.59 & -0.019800 \\
\bottomrule
\end{tabular}
\end{table}

The most negative energy achieved for \(\mathrm{H_2}\) was \(-1.13730605\,\mathrm{Ha}\), with a difference from the exact FCI energy \(-1.117506\,\mathrm{Ha}\). Step counts of 10 provide an optimal trade-off between convergence speed and computational cost.

\subsubsection{Molecule \(\mathrm{LiH}\)}

\begin{table}[H]
\centering
\caption{Final energy and optimization time for \(\mathrm{LiH}\) using Momentum optimizer with varying steps.}
\begin{tabular}{lccc}
\toprule
\textbf{Steps} & \textbf{Final Energy (Ha)} & \textbf{Total Time (s)} & \textbf{Difference from FCI (Ha)} \\
\midrule
5 & \(-7.88079149\) & \(\mathbf{8328.30}\) & -0.017409 \\
10 & \(-7.88118152\) & 13365.38 & \(\mathbf{-0.017799}\) \\
15 & \(-7.88226478\) & 18384.78 & -0.018883 \\
20 & \(-7.84260017\) & 23309.87 & 0.020782 \\
25 & \(-7.88176386\) & 28433.90 & -0.018382 \\
30 & \(-7.88222895\) & 33421.15 & -0.018847 \\
\bottomrule
\end{tabular}
\end{table}

For \(\mathrm{LiH}\), the energy \(-7.88226478\,\mathrm{Ha}\) achieved in 15 steps shows the best accuracy compared to the FCI value \(-7.863382\,\mathrm{Ha}\), but the computation time is significantly lower for 5 steps (\(8328.30\,\mathrm{s}\)).

\subsubsection{Conclusions}

The results indicate that:
- For \(\mathrm{H_2}\), 10 steps balance speed and energy convergence.
- For \(\mathrm{LiH}\), 15 steps provide the most accurate energy, but fewer steps (5 or 10) can save substantial time with acceptable accuracy trade-offs.



\section{Limitations}
In the development of this project, several limitations have been identified that could have influenced the results obtained. It is essential to recognize and understand these limitations to properly interpret the conclusions and guide future research efforts.

\subsection{Availability of a Quantum Computer}
One of the most evident limitations has been the lack of access to a real quantum computer to validate the developed implementation. The simulation was conducted on classical hardware, which means it was not possible to directly verify how the simulator behaves on an actual quantum platform. This restriction limits the ability to evaluate the performance and efficiency of the implementation in a real quantum environment, where factors such as quantum coherence and entanglement could significantly influence the results.

\subsection{No Implementation of Quantum Errors}
Another significant limitation has been the absence of quantum error modeling in the simulation. In real quantum computers, errors arising from decoherence, imperfect quantum gates, and environmental noise are inevitable and can drastically affect the accuracy of simulations. By not incorporating these errors into the model, the simulation lacks the robustness required to faithfully reflect the operating conditions of a real quantum system. Including errors could have provided greater authenticity to the results, allowing for an evaluation of the simulator's resilience to the inherent imperfections of quantum computing.





%%%% BUDGET %%%%
% Budget chapter. Please replace "budget.tex" entirely with your own content written in the desired language.
%\input{budget}


%%%% SUSTAINABILITY REPORT %%%%    <<<<<<<<<<<<<<<< NEW !!
% Sustainability chapter. Please replace "sustainability.tex" entirely with your own content written in the desired language.
\chapter{Sustainability Analysis and Ethical Implications}

\section{Sustainability Matrix}
This document presents an overview of the project’s sustainability by examining three key perspectives—environmental, economic, and social—across the distinct phases of development, execution, and potential risks or limitations.

\subsection{Environmental Perspective}

\subsubsection{Development}

The development phase of the project required approximately 540 hours of work. To carry out this work, I used a Mac laptop with an average power consumption of 30\,W and a high-performance computing (HPC) server provided by a Viennese research center. This server, equipped with 16 CPU threads, operates at around 200\,W under heavy workloads. Combining these resources allowed me to optimize performance while avoiding additional hardware acquisition or excessive energy consumption.

Quantum simulations were managed using the open-source \emph{PennyLane} framework, which, in this case, did not require GPU acceleration, thus keeping power demands moderate. Additionally, I minimized environmental impact during commuting by traveling approximately 3.5\,km per day via metro over 120 days. Factoring in all these contributions, the total emissions for the project are estimated to be 32.73\,kg of CO\(_2\). While I did not explicitly follow a circular economy model, I prioritized reusing existing infrastructure and avoided purchasing new devices, which further reduced potential waste.

About de materials and resources origins, no fresh hardware purchases were made. I relied on the Viennese HPC center’s existing infrastructure, which follows European directives on responsible energy usage and publishes annual reports discussing sustainability practices. \emph{PennyLane} itself is developed by a community that strives for ethical and transparent technology, so both the HPC resources and the software align with these principles.

The emissions calculations are detailed below:
\begin{ProjectStructure}
    \textbf{Laptop emissions:}
    \[
    E_{\text{l}} = P_{\text{laptop}} \cdot t_{\text{laptop}} \cdot \text{EF}_{\text{electricity}} = 0.03 \, \text{kW} \cdot 540 \, \text{h} \cdot 0.25 \, \text{kg CO}_2/\text{kWh} = 4.05 \, \text{kg CO}_2 
    \]
    \textbf{Server emissions:}
    \[
    E_{\text{s}} = P_{\text{server}} \cdot t_{\text{server}} \cdot \text{EF}_{\text{electricity}} = 0.2 \, \text{kW} \cdot 540 \, \text{h} \cdot 0.25 \, \text{kg CO}_2/\text{kWh} = 27.0 \, \text{kg CO}_2
    \]
    \textbf{Transport emissions:}
    \[
    E_{\text{t}} = d_{\text{metro}} \cdot n_{\text{days}} \cdot \text{EF}_{\text{metro}} = 3.5 \, \text{km} \cdot 120 \, \text{days} \cdot 0.014 \, \text{kg CO}_2/\text{km} = 1.68 \, \text{kg CO}_2
    \]
    \textbf{Total emissions:}
    \[
    E_{\text{total}} = E_{\text{l}} + E_{\text{s}} + E_{\text{t}} = 4.05 \, \text{kg CO}_2 + 27.0 \, \text{kg CO}_2 + 1.68 \, \text{kg CO}_2 = 32.73 \, \text{kg CO}_2
    \]
\end{ProjectStructure}

\subsubsection{Execution}
Once the project moves beyond its initial development and into practical use, the primary ongoing resource becomes CPU time, both on the laptop and at the Viennese HPC center. I anticipate an annual energy consumption in the range of 500\,kWh for active simulations, leading to approximately 75\,kg of CO\(_2\) if the current power mix remains the same. However, because the HPC center integrates some lower-carbon energy sources, there is hope that the carbon intensity of these computations could decrease over time.

One of the key advantages of leveraging quantum approaches is the possibility of reducing the energy used in large-scale calculations. Compared to purely classical methods, certain quantum-inspired algorithms—especially when optimized through \emph{PennyLane}—can cut CPU time, which in turn saves up to an estimated 300\,kWh per year. This reduction lowers both energy costs and carbon emissions.

At the conclusion of the project, as it is primarily a digital endeavor, no physical waste is generated. Upon completion, the data can be either archived or securely deleted, requiring minimal energy consumption and resulting in an almost negligible environmental footprint. Relying on renewable energy within the HPC center, continuing to refine algorithms, and scaling down usage when computational tasks are not urgent could all reduce the overall footprint further.

Emissions calculations:
\begin{ProjectStructure}
    \textbf{Annual energy consumption emissions:}
    \[
    E_{\text{annual}} = E_{\text{cpu}} \cdot \text{EF}_{\text{electricity}} = 500 \, \text{kWh} \cdot 0.15 \, \text{kg CO}_2/\text{kWh} = 75 \, \text{kg CO}_2
    \]
    \textbf{Potential savings from quantum optimization:}
    \[
    E_{\text{saved}} = \Delta E_{\text{cpu}} \cdot \text{EF}_{\text{electricity}} = 300 \, \text{kWh} \cdot 0.15 \, \text{kg CO}_2/\text{kWh} = 45 \, \text{kg CO}_2
    \]
    \textbf{Net emissions with optimization:}
    \[
    E_{\text{net}} = E_{\text{annual}} - E_{\text{saved}} = 75 \, \text{kg CO}_2 - 45 \, \text{kg CO}_2 = 30 \, \text{kg CO}_2
    \]
\end{ProjectStructure} 
\subsubsection{Risks and Limitations}
Unplanned expansions in testing or running extended, unoptimized simulations on the HPC server could quickly increase energy consumption. Without proper scheduling or oversight, the footprint could grow significantly.
If I repeated this work, I would consider other cloud-based quantum services that have carbon-neutral certifications or advanced versions of \emph{PennyLane} with improved efficiency. Tight scheduling of simulations to off-peak energy hours could also help.
Key figures rely on average emission factors, and the power mix of the HPC center can change. Precise measurements for every hour of server usage are challenging, which introduces uncertainty into the calculations.


\subsection{Economic Perspective}

\subsubsection{Development}
The majority of expenses during the development phase came from labor. I devoted around 540 hours to the project, valued at 25\,\euro/h, for a total of 13{,}500\,\euro. Electricity costs for the workstation added approximately 20\,\euro, assuming an average price of 0.18\,\euro/kWh for the roughly 108\,kWh consumed. Because I relied on the Viennese HPC center under a research arrangement, access to high-performance infrastructure incurred no direct additional fees. The open-source nature of \emph{PennyLane} likewise avoided software licensing costs and minimized overall expenses.

\subsubsection{Execution}
If the project remains active and consumes around 500\,kWh annually, it would cost an additional 90\,\euro\ per year (at 0.18\,\euro/kWh). Improved energy prices or algorithmic refinements could further lower this expense.

Most maintenance revolves around software updates, bug fixes, and code improvements, which involve minimal monetary outlays. The Viennese HPC center supports open-source frameworks, and \emph{PennyLane} receives community-backed updates at no extra cost.

Since the project is entirely digital, shutting it down incurs negligible expense. Data can be archived or securely deleted, and no physical equipment needs special handling.

Any quantum simulation tools or modules developed here could assist future research, decreasing start-up costs and encouraging interdisciplinary collaboration. The open approach ensures that others can build on these methods freely.


\subsubsection{Risks and Limitations}
If energy prices were to soar or the HPC center discontinued free access, the project’s costs could skyrocket, potentially hindering further progress.

Estimates rely on current prices and the stable availability of HPC resources. Rapid market changes, especially in energy, might invalidate long-term cost projections.

\subsection{Social Perspective}

\subsubsection{Development}
Throughout the development stage, I emphasized ethically and socially responsible practices. This included using clear, inclusive language in all documentation and making the core results openly accessible. Given that quantum computing is a pioneering area with potential global impact, I aimed to set a collaborative tone and steer clear of any design choices that might discriminate against particular user groups.

The Viennese HPC center publishes regular reports outlining its commitment to responsible energy use and fair resource distribution. Meanwhile, the open-source community around \emph{PennyLane} encourages knowledge sharing and keeps barriers to entry low, which helps cultivate a more equitable research environment.

\subsubsection{Execution}
As the project transitions into continuous operation, a diverse audience of researchers in fields such as computational chemistry, quantum physics, and beyond may benefit from these optimized simulations. Users anywhere in the world can deploy the code, provided they have sufficient HPC access, and the open documentation helps ensure they understand and can modify the methodology. However, disparities in HPC availability could widen the gap between institutions that can leverage quantum resources effectively and those that cannot.

In addressing the original challenge of achieving faster, more efficient simulations, the solution outlined here has proven successful. By tapping into HPC resources under carefully designed algorithms, the project shows how quantum-inspired frameworks, like \emph{PennyLane}, can reduce computational overhead without requiring proprietary software.

\subsubsection{Risks and Limitations}
A key concern relates to uneven access. If HPC centers are concentrated in a few regions, researchers in other parts of the world may struggle to replicate results or remain competitive. Another vulnerability surfaces if \emph{PennyLane} or HPC usage policies become more restrictive, potentially locking in users who have shaped their workflows around these platforms. Finally, it is worth noting that this social analysis assumes relatively stable conditions that may not always apply to low-income or remote communities, where internet connectivity and funding are limited.

\section{Ethical Implications}
The project aligns with the university’s Code of Ethics by promoting open research and striving to use energy responsibly. Quantum computing carries the potential to transform a wide range of scientific activities, so I made sure to prioritize transparency, collaboration, and accessibility. Each phase of the work underscores an intention to minimize harmful applications, support inclusive participation, and encourage responsible innovation that extends benefits to the broader community.

\section{Relation to the Sustainable Development Goals (SDGs)}
\begin{itemize}
    \item \textbf{Goal 9 (Industry, Innovation, and Infrastructure):}  
    By taking advantage of high-performance computing and refining quantum simulation techniques, the project fosters advances in scientific research and knowledge dissemination.
    \item \textbf{Goal 13 (Climate Action):}  
    Through careful algorithmic design and reduced computational overhead, the project mitigates carbon emissions and emphasizes strategies to leverage clean energy sources, aligning with global climate objectives.
\end{itemize}


%%%% CONCLUSIONS AND FUTURE WORK %%%%
% Conclusions chapter. Please replace "conclusions.tex" entirely with your own content written in the desired language.
\chapter{Conclusions and Future Work}

\section{Conclusions}

In this project, we have successfully implemented a quantum simulator focused on simulating various molecular systems using the VQE algorithm. The results demonstrate the ability to optimize the processes for H$_2$, LiH, and H$_2$O, molecules of different sizes and complexities. Additionally, the efficacy of the UCCSD ansatz has been validated, achieving superior results compared to other ansätze. An unexpected outcome, which significantly enhances the value of our project, was demonstrating that, for our implementation, the \textit{autograd} interface outperforms JAX.

The contributions of this work to the field include the modular design of quantum simulation software, enabling extensibility and flexibility for future algorithms and optimization strategies. Furthermore, the introduction of an adaptive methodology for selecting different ansätze and optimizers has made the optimization process more efficient and faster.

Consequently, the objectives of our project have been successfully achieved. These objectives include the development of an adaptive ansatz capable of selecting the operators with the greatest impact on the system, the integration of a hybrid optimization approach that adjusts both the ansatz parameters and nuclear positions, and endows our simulations with greater robustness and precision. Finally, the design of a modular architecture facilitates the inclusion of new types of ansätze and optimizers, enabling the comparison of different optimization strategies for future research.

\section{Future Directions}
In this project, the initial questions posed have been successfully addressed. However, after the completion of the project, new frontiers and research lines have emerged, offering exciting opportunities for future work.

During the development of the project, the hybrid optimization methodology could be further refined by adding to the optimization process, thereby achieving greater efficiency in convergence. We realized that another viable option to enhance simulation convergence would be to implement an adaptive optimization system, which dynamically adjusts the number of optimizations of the ansatz parameters based on the simulation’s progress. This would result in greater efficiency in the simulation’s convergence. Moreover, it was observed that automating the configuration of optimizers would save considerable time during setup.

As a final future endeavor, to improve the simulator's consistency, it would be valuable to investigate its performance on NISQ systems. This would enable analyzing whether implementing a noise model affects the simulation’s convergence, thus validating the robustness of the simulator in real-world quantum systems.

Another relevant research avenue would involve integrating error mitigation techniques, such as zero-noise extrapolation, to enhance the accuracy of simulations on noisy quantum devices. This would add greater realism to the simulations and make the simulator more practical for applications on current quantum hardware.

Finally, given the simulator’s modular structure, it could be extended to other fields beyond chemistry, such as material science and condensed matter physics. Exploring these areas would broaden its impact and leverage its flexibility for a wide range of scientific applications.


%%%% BIBLIOGRAPHY %%%%
\nocite{*}               % Forces all entries to be printed, even if not cited
% Bibliography intro
% Supress this macro (or modify its contents to suit your needs).
\defbibnote{bib-intro}{%
\ifcase\doclanguage\or
  El sistema \textit{biblatex} simplifica la gestió de la bibliografia en treballs científics, proporcionant automatització i personalització en el format de les citacions. Això permet a l'autor del document enfocar-se en el contingut sense haver de preocupar-se per l'estil de les referències, estalviant temps i reduint errors.
  
  La base de dades de referències bibliogràfiques és al fitxer «TFG.bib» i és allà on heu d'afegir les vostres referències. Consulteu el manual del \texttt{biblatex}, secció «Database Guide», per conèixer els tipus de referències i camps disponibles.
  
  Podeu modificar (o suprimir) aquesta nota editant la macro \texttt{\textbackslash defbibnote} al fitxer «TFG.tex».
  \par\hfil\rule[3pt]{.5\textwidth}{0.4pt}\hfil\par\or
  El sistema \textit{biblatex} simplifica la gestión de la bibliografía en trabajos científicos, proporcionando automatización y personalización en el formato de las citas. Esto permite al autor del documento enfocarse en el contenido sin tener que preocuparse por el estilo de las referencias, ahorrando tiempo y reduciendo errores.
  
  La base de datos de referencias bibliográficas está en el archivo «TFG.bib» y es allí donde se deben añadir vuestras referencias. Consultar el manual de \texttt{biblatex}, sección «Database Guide», para conocer los tipos de referencias y campos disponibles.
  
  Podéis modificar (o suprimir) esta nota editando la macro \texttt{\textbackslash defbibnote} en el archivo «TFG.tex».
  \par\hfil\rule[3pt]{.5\textwidth}{0.4pt}\hfil\par\else
  The \textit{biblatex} system simplifies the management of the bibliography in scientific works, providing automation and customization in the format of citations. This allows the document's author to focus on the content without having to worry about the style of the references, saving time and reducing errors.
  
  The bibliographic references database is in the file “TFG.bib”, and this is where you should add your references. Consult the \texttt{biblatex} manual, section “Database Guide”, to learn about the available types of references and fields.
  
  You can modify (or delete) this note by editing the \texttt{\textbackslash defbibnote} macro in the file “TFG.tex”.
  \par\hfil\rule[3pt]{.5\textwidth}{0.4pt}\hfil\par\fi%
}
\printbibliography[heading=bibintoc,prenote={bib-intro}]


%%%% ANNEXES %%%%
% All chapters AFTER the \appendix command wil be considered appendices and numbered by letter
\appendix
\chapter{Logic Gates}
\label{appendices:LogicGates}
\section{Simple Logic Gates}
Below are detailed the simple logic gates essential for constructing more complex quantum algorithms:

\paragraph{X Gate (Pauli-X)}

The \textbf{Pauli-X} gate is the quantum analog of the classical NOT gate. It performs a bit flip on the qubit, transforming the state $\ket{x}$ into $\ket{\neg x}$.

\vspace{0.5em}
\begin{minipage}{\textwidth}
    \begin{minipage}[t]{0.45\textwidth}
        \centering
        \textbf{Representative Matrix}\\[0.5em]
        \[
        X = 
        \begin{pmatrix}
        0 & 1 \\
        1 & 0 \\
        \end{pmatrix}
        \]
    \end{minipage}
    \hfill
    \begin{minipage}[t]{0.45\textwidth}
        \centering
        \textbf{Effect on Basis States}\\[0.5em]
        \begin{itemize}
            \item $X\ket{0} = \ket{1}$
            \item $X\ket{1} = \ket{0}$
        \end{itemize}
    \end{minipage}
\end{minipage}

\paragraph{Y Gate (Pauli-Y)}

The \textbf{Pauli-Y} gate performs a rotation of $\pi$ around the $y$-axis. It transforms the state $\ket{x}$ into $i(-1)^x\ket{\neg x}$.

\vspace{0.5em}

\begin{minipage}{\textwidth}
    \begin{minipage}[t]{0.45\textwidth}
        \centering
        \textbf{Representative Matrix}\\[0.5em]
        \[
        Y = 
        \begin{pmatrix}
        0 & -i \\
        i & 0 \\
        \end{pmatrix}
        \]
    \end{minipage}
    \hfill
    \begin{minipage}[t]{0.45\textwidth}
        \centering
        \textbf{Effect on Basis States}\\[0.5em]
        \begin{itemize}
            \item $Y\ket{0} = i\ket{1}$
            \item $Y\ket{1} = -i\ket{0}$
        \end{itemize}
    \end{minipage}
\end{minipage}

\paragraph{Z Gate (Pauli-Z)}

The \textbf{Pauli-Z} gate is known as the phase inversion gate. It transforms the state $\ket{x}$ into $(-1)^x\ket{x}$.

\vspace{0.5em}
\begin{minipage}{\textwidth}
    \begin{minipage}[t]{0.45\textwidth}
        \centering
        \textbf{Representative Matrix}\\[0.5em]
        \[
        Z = 
        \begin{pmatrix}
        1 & 0 \\
        0 & -1 \\
        \end{pmatrix}
        \]
    \end{minipage}
    \hfill
    \begin{minipage}[t]{0.45\textwidth}
        \centering
        \textbf{Effect on Basis States}\\[0.5em]
        \begin{itemize}
            \item $Z\ket{0} = \ket{0}$
            \item $Z\ket{1} = -\ket{1}$
        \end{itemize}
    \end{minipage}
\end{minipage}


\paragraph{Hadamard Gate (H)}

The \textbf{Hadamard} gate creates an equal superposition of the computational basis states. It transforms the state $\ket{x}$ into $\dfrac{1}{\sqrt{2}} (\ket{0} + (-1)^x \ket{1})$.

\vspace{0.5em}
\begin{minipage}{\textwidth}
    \begin{minipage}[t]{0.45\textwidth}
        \centering
        \textbf{Representative Matrix}\\[0.5em]
        \[
        H = \dfrac{1}{\sqrt{2}}
        \begin{pmatrix}
        1 & 1 \\
        1 & -1 \\
        \end{pmatrix}
        \]
    \end{minipage}
    \hfill
    \begin{minipage}[t]{0.45\textwidth}
        \centering
        \textbf{Effect on Basis States}\\[0.5em]
        \begin{itemize}
            \item $H\ket{0} = \dfrac{1}{\sqrt{2}} (\ket{0} + \ket{1}) = \ket{+}$
            \item $H\ket{1} = \dfrac{1}{\sqrt{2}} (\ket{0} - \ket{1}) = \ket{-}$
        \end{itemize}
    \end{minipage}
\end{minipage}

\section{Multi-Qubit Logic Gates}

\paragraph{Controlled-NOT Gate (CNOT)}

The \textbf{CNOT} or \textbf{Controlled-X} gate is a two-qubit gate that flips the second qubit (target) if and only if the first qubit (control) is in the state $\ket{1}$. It transforms the state $\ket{x, y}$ into $\ket{x, x \oplus y}$, where $\oplus$ denotes the XOR operation.

\vspace{0.5em}
\begin{minipage}{\textwidth}
    \begin{minipage}[t]{0.45\textwidth}
        \centering
        \textbf{Representative Matrix}\\[0.5em]
        \[
        \text{CNOT} = 
        \begin{pmatrix}
        1 & 0 & 0 & 0 \\
        0 & 1 & 0 & 0 \\
        0 & 0 & 0 & 1 \\
        0 & 0 & 1 & 0 \\
        \end{pmatrix}
        \]
    \end{minipage}
    \hfill
    \begin{minipage}[t]{0.45\textwidth}
        \centering
        \textbf{Effect on Basis States}\\[0.5em]
        \begin{itemize}
            \item $\text{CNOT}\ket{00} = \ket{00}$
            \item $\text{CNOT}\ket{01} = \ket{01}$
            \item $\text{CNOT}\ket{10} = \ket{11}$
            \item $\text{CNOT}\ket{11} = \ket{10}$
        \end{itemize}
    \end{minipage}
\end{minipage}

\paragraph{Single Excitation Gate (\textit{SingleExcitation})}

This gate performs a rotation in the two-dimensional subspace $\{\ket{01}, \ket{10}\}$. It transforms the state $\ket{10}$ into $\cos\left( \dfrac{\phi}{2} \right) \ket{10} - \sin\left( \dfrac{\phi}{2} \right) \ket{01}$.

\vspace{0.5em}
\begin{minipage}{\textwidth}
    \begin{minipage}[t]{0.45\textwidth}
        \centering
        \textbf{Representative Matrix}\\[0.5em]
        \[
        U(\phi) = \begin{pmatrix}
        1 & 0 & 0 & 0 \\
        0 & \cos\left( \dfrac{\phi}{2} \right) & -\sin\left( \dfrac{\phi}{2} \right) & 0 \\
        0 & \sin\left( \dfrac{\phi}{2} \right) & \cos\left( \dfrac{\phi}{2} \right) & 0 \\
        0 & 0 & 0 & 1 \\
        \end{pmatrix}
        \]
    \end{minipage}
    \hfill
    \begin{minipage}[t]{0.45\textwidth}
        \centering
        \textbf{Effect on Basis States}\\[0.5em]
        It affects the subspace $\{\ket{01}, \ket{10}\}$, performing a rotation parameterized by $\phi$.
    \end{minipage}
\end{minipage}

\paragraph{Double Excitation Gate (\textit{DoubleExcitation})}

This gate performs a rotation in the subspace of states $\{\ket{0011}, \ket{1100}\}$. It specifically affects these states, leaving the others unchanged.

\vspace{0.5em}
\begin{minipage}{\textwidth}
    \begin{minipage}[t]{0.45\textwidth}
        \centering
        \textbf{Representative Matrix}\\[0.5em]
        \[
        U(\phi) = \begin{pmatrix}
        I_{12} & 0 & 0 \\
        0 & \begin{pmatrix}
        \cos\left( \dfrac{\phi}{2} \right) & -\sin\left( \dfrac{\phi}{2} \right) \\
        \sin\left( \dfrac{\phi}{2} \right) & \cos\left( \dfrac{\phi}{2} \right) \\
        \end{pmatrix} & 0 \\
        0 & 0 & I_2 \\
        \end{pmatrix}
        \]
    \end{minipage}
    \hfill
    \begin{minipage}[t]{0.45\textwidth}
        \centering
        \textbf{Effect on Basis States}\\[0.5em]
        It performs a rotation parameterized by $\phi$ in the subspace $\{\ket{0011}, \ket{1100}\}$.
    \end{minipage}
\end{minipage}

These gates are implemented in PennyLane as \texttt{qml.SingleExcitation} and \linebreak\texttt{qml.DoubleExcitation}, and are essential in quantum chemistry algorithms such as the \textit{Unitary Coupled-Cluster Singles and Doubles} (UCCSD).


%%%% FINAL CLOSING BLANK PAGE %%%%
\clearpage
\thispagestyle{empty}
~

\end{document}
